
\part{시뮬레이션 알고리즘}
\label{part:algorithm}

\chapter{항목별 데이터 변환 방법}
\label{chap:simulator_geometry}

\section{시뮬레이션 설정}
\label{section:global-model-assumptions}

모든 모델에는 동일한 EnergyPlus 버전, 시간간격, 계산기간, 크기산정기간 및 대류열전달 설정을 적용한다.

\begin{defaulttable}
  \caption{모든 모델에 공통으로 적용되는 EnergyPlus 계산 설정}
  \label{tab:global-energyplus-settings}
  \small
  \begin{tabularx}{\textwidth}{p{4cm}p{4.1cm}X}
    \toprule
    항목 & 적용값 & 생성 내용 \\ \midrule
    EnergyPlus 버전 & 24.2.0 & 해당 버전의 IDD와 객체 필드를 사용 \\
    계산 시간간격 & 시간당 6회 & 10분 간격 \\
    연간 계산기간 & 1월 1일--12월 31일 & 365일 계산 \\
    달력 연도 & 2026년 & 요일 기반 스케줄의 날짜·요일 연결에 사용 \\
    난방 크기산정 기간 & 기상파일의 1월 1일--31일 & \texttt{SizingPeriod:WeatherFileDays} 생성 \\
    냉방 크기산정 기간 & 기상파일의 8월 1일--31일 & \texttt{SizingPeriod:WeatherFileDays} 생성 \\
    실내·실외 대류모델 & TARP & 존 안쪽 및 바깥쪽 표면에 적용 \\
    존 난방 크기산정 계수 & 1.25 & \texttt{Sizing:Zone}에 적용 \\
    존 냉방 크기산정 계수 & 1.15 & \texttt{Sizing:Zone}에 적용 \\ \bottomrule
  \end{tabularx}
\end{defaulttable}

용량과 유량을 생략한 설비 필드에는 가능한 경우 \texttt{Autosize}를 입력한다. 크기산정 계산은 위 표의 1월 및 8월 기상파일 기간과 각 설비 절에 제시된 설계온도·설계온도차를 사용한다.


\section{형상}

\subsection{면: Detailed Geometry (DG) 와 Simplified Geometry (SG)}

건물 에너지 시뮬레이션 도구의 형상 입력 방식은 크게 Detailed Geometry (DG)와 Simplified Geometry (SG)로 구분할 수 있다\parencref{fig:geometry_types}. Detailed Geometry 방식은 실의 각 면들을 구성하는 점들의 좌푯값을 규칙에 맞게 상세하게 입력하여 정의하는 방식으로, 실재하는 건물의 3D 형상을 그대로 구현하는 방식이다. EnergyPlus의 기본적인 형상 정의 방식으로 DesignBuilder, OpenStudio, ClimateStudio 등 EnergyPlus를 기반으로 하는 일반적인 건물 에너지 시뮬레이션 도구에서 사용되고 있다.

\begin{defaultfigure}
  
  \begin{subfigure}[b]{0.45\textwidth}
    \centering
    \includegraphics[scale=0.6]{형상_DG.png}
    \caption{Detailed Geometry (DG)}
  \end{subfigure}
  \hfill
  \begin{subfigure}[b]{0.45\textwidth}
    \centering
    \includegraphics[scale=0.6]{형상_SG.png}
    \caption{Simplified Geometry (SG)}
  \end{subfigure}
  
  \caption{형상 입력 방식}
  \label{fig:geometry_types}

\end{defaultfigure}

반면 Simplified Geometry 방식은 실의 각 면들을 그 면적과 향(방위각) 만으로 정의하는 방식으로, 열평형 방정식 기반으로 풀어지는 건물 에너지 시뮬레이션에서 핵심이 되는 정보만 입력을 받는 개념으로 볼 수 있다. 좌표값의 계산, 3D 모델링 도구 등 없이 면을 나열할 수 있다는 점에서 높은 편의성을 가지며, ISO 13790 등 준정적 시뮬레이션 도구들에서 흔히 사용되며, EnergyPlus를 기반으로 하더라도 도구의 사용 편의성, 외부 프로그램 비의존성을 형상의 시각화, 상세한 일사에너지 계산보다 우선시하는 경우 종종 사용된다 \cite{cbecc}.

\simulator 는 SG 방식으로 형상을 입력받아, 이를 EnergyPlus 시뮬레이션 가능한 DG 정보로 변환하는 방식을 취하고 있다.EnergyPlus에서 형상은, 그 면적과 향 및 경계조건이 동일하다면, 상호 그림자에 의한 일사 차폐 효과를 무시하는 한 열적으로 동일하다. 따라서, 면적과 향으로 주어진 각 면들은 모두 원점을 지나는 임의의 평면으로 생성되며, EnergyPlus에서 요구하는 인접 존 처리 규칙을 만족시키기 위해 두 존 사이의 벽체들은 모두 복제(clone 및 flip)되며, 이 과정에서 입력된 구조체도 복제 및 뒤집어져 적용된다 \parencref{fig:sg-to-dg-vertex-extraction}. 

\begin{defaultfigure}
  \includegraphics{sg-to-dg-vertex-extraction.png}
  \caption{Simplified Geometry 정보의 Detailed Geometry 변환 과정}
  \label{fig:sg-to-dg-vertex-extraction}
\end{defaultfigure}

수평면인 바닥과 천장은 입력 면적 $A_h$를 갖는 정사각형으로 생성한다. 한 변의 길이 $L_h$는 다음과 같다.

\begin{equation}
  L_h = \sqrt{A_h}.
  \label{eq:equivalent-horizontal-surface}
\end{equation}

바닥은 $z=0$, 천장은 해당 존의 천장고 $H_z$에 배치한다. 벽은 높이를 존 천장고와 같게 하고, 폭 $W_w$를 다음과 같이 정한 직사각형으로 생성한다.

\begin{equation}
  W_w = \frac{A_w}{H_z},
  \label{eq:equivalent-wall-width}
\end{equation}

여기서 $A_w$는 개구부 면적을 포함한 벽 전체 면적이다. 외기 벽은 입력 방위각을 이용해 법선 방향을 정한다. 바닥, 천장 및 외기에 노출되지 않은 벽의 좌표상 위치와 향은 실제 도면 위치를 재현하지 않으며, 열적 연결을 구성하기 위한 임의의 위치에 생성된다. 건물의 \inlinecode{north_axis}는 생성된 건물 전체에 EnergyPlus의 건물 회전값으로 적용된다.

이렇게 EnergyPlus에서 구현된 DG 정보는, 원래 건물의 3D 형상을 복원한 것이 아닌, 이와 열적으로 동등한 3D 면들의 집합체이다 \parencref{fig:transformed-sg-example}. 상호 그림자에 의해 일사가 차폐되는 것을 막기 위해 상호간 그림자는 완전히 무시되는 조건으로 시뮬레이션을 실행한다.

\begin{defaultfigure}
  \includegraphics{transformed-sg-example.png}
  \caption{Detailed Geometry로 변환된 Simplified Geometry의 예시}
  \label{fig:transformed-sg-example}
\end{defaultfigure}


\subsection{개구부}

각 개구부는 소속면의 로컬 좌표 $(0.001,0.001)$에서 시작하는 정사각형으로 생성한다. 입력 면적을 $A_f$라 하면 EnergyPlus 객체에 전달하는 가로와 세로는 모두 $0.999\sqrt{A_f}$이다. 동일 면에 여러 개구부가 있으면 각 객체에 같은 시작좌표를 사용한다.

\texttt{window}와 \texttt{glassdoor}는 투명 개구부로 분류되어 \texttt{Window} 객체로 변환되고, \texttt{door}는 \texttt{Door} 객체로 변환된다. 존 간 면에서는 각각 \texttt{Window:Interzone} 및 \texttt{Door:Interzone}을 사용하며, 상대면의 개구부 목록을 같은 순서로 연결한다. 외기, 지면 및 단열 경계에서는 일반 \texttt{Window} 또는 \texttt{Door} 객체를 만들고 소속 건물면의 이름을 모체면으로 기록한다. EnergyPlus는 모체면 경계조건과 개구부 객체의 조합을 입력검사 단계에서 검증한다.

\section{재료 및 구조체}

사용자가 재료층을 명시한 구조체는 입력된 열전도율, 밀도, 비열 및 두께로 EnergyPlus의 \texttt{Material}과 \texttt{Construction} 객체를 생성한다. 입력하지 않은 표면 광학속성은 재료 거칠기 \texttt{Rough}, 장파복사 흡수율 0.9, 태양열 흡수율 0.7 및 가시광선 흡수율 0.7을 사용한다.

\subsection{\texttt{unknown} 구조체의 해소}

\inlinecode{construction_id}가 \inlinecode{unknown}인 면은 단열 경계 또는 열전달이 없는 면으로 간주하지 않는다. \simulator 는 EnergyPlus 입력파일을 생성하기 직전에 각 \inlinecode{unknown} 면을 실제 재료층을 갖는 구조체로 치환한다.

\subsubsection{면 유형에 따른 구조체 선택}

\paragraph{내벽}
면 유형이 \inlinecode{wall}이고 경계조건이 \inlinecode{zone}이면 법규
열관류율 DB를 조회하지 않고, 다음의 대칭 경량 내벽을 적용한다.

\[
  \text{석고보드 }12.5\ \mathrm{mm}
  \;-\;
  \text{글라스울 }50.0\ \mathrm{mm}
  \;-\;
  \text{석고보드 }12.5\ \mathrm{mm}
\]

\paragraph{그 외}
내벽을 제외한 \inlinecode{unknown} 면은 법규 열관류율 DB를 조회한다. 면 유형과 경계조건은 \cref{tbl:unknown-surface-part-mapping}과 같이 DB의 부위명으로 변환된다.

\begin{defaulttable}
  \caption{\texttt{unknown} 면의 법규 DB 부위 결정}
  \label{tbl:unknown-surface-part-mapping}
  {\small
  \begin{tabularx}{\textwidth}{lllX}
    \toprule
    면 유형 & 경계조건 & DB 부위 & 처리 내용 \\ \midrule
    \texttt{wall}
      & \texttt{zone}
      & 해당 없음
      & 석고보드--글라스울--석고보드 기본 내벽 적용 \\

    \texttt{wall}
      & \texttt{outdoors}, \texttt{ground}, \texttt{adiabatic}
      & 외벽
      & 외벽 열관류율 조회 \\

    \texttt{ceiling}
      & \texttt{outdoors}
      & 최상층 지붕
      & 최상층 지붕 열관류율 조회 \\

    \texttt{ceiling}
      & \texttt{zone}, \texttt{ground}, \texttt{adiabatic}
      & 바닥난방이 아닌 층간바닥
      & 비난방 층간바닥 열관류율 조회 \\

    \texttt{floor}
      & \texttt{outdoors}, \texttt{ground}
      & 바닥난방이 아닌 최하층 바닥
      & 비난방 최하층 바닥 열관류율 조회 \\

    \texttt{floor}
      & \texttt{zone}, \texttt{adiabatic}
      & 바닥난방이 아닌 층간바닥
      & 비난방 층간바닥 열관류율 조회 \\
    \bottomrule
  \end{tabularx}
  }
\end{defaulttable}

법규 DB의 조회 키는 다섯 항목(시행일자, 부위, 외기조건, 용도, 기후지역)으로 구성되며, 이 중 부위는 위 표에서 결정되고, 나머지 네 항목은 다음과 같이 결정된다.

\begin{itemize}
  \item 용도: \texttt{공동주택 외}
  \item 외기조건: \texttt{외기 직접}
  \item 바닥 구분: \texttt{바닥난방이 아닌}
  \item 기후지역: \texttt{중부1}, \texttt{중부2}, \texttt{남부}, \texttt{제주도} 중 주소와 허가일에 따라 결정된 값
\end{itemize}

\subsubsection{시행일자의 선택}

건물 허가일보다 이전인 시행일 중 가장 최근 시행일을 선택한다.

\subsubsection{내장 법규 열관류율 DB}

아래 표의 모든 값은 목표 열관류율이며 단위는
$\mathrm{W/(m^2K)}$이다.

\begin{defaulttable}
  \caption{\texttt{unknown} 외벽에 적용되는 열관류율
    $[\mathrm{W/(m^2K)}]$}
  \label{tbl:unknown-db-wall}
  {\small
  \renewcommand{\arraystretch}{1.05}
  \begin{tabular}{lrrrr}
    \toprule
    시행일자 & 중부1 & 중부2 & 남부 & 제주도 \\ \midrule
    2018-09-01 & 0.17 & 0.24 & 0.32 & 0.41 \\
    2016-01-01 & 0.26 & 0.26 & 0.32 & 0.43 \\
    2013-10-01 & 0.27 & 0.27 & 0.34 & 0.44 \\
    2013-09-01 & 0.27 & 0.27 & 0.34 & 0.44 \\
    2010-11-05 & 0.36 & 0.36 & 0.45 & 0.58 \\
    2008-07-10 & 0.47 & 0.47 & 0.58 & 0.76 \\
    2002-08-31 & 0.47 & 0.47 & 0.58 & 0.76 \\
    2001-01-17 & 0.47 & 0.47 & 0.58 & 0.76 \\
    1996-02-09 & 0.58 & 0.58 & 0.76 & 1.16 \\
    1992-06-01 & 0.58 & 0.58 & 0.76 & 1.16 \\
    1987-07-21 & 0.58 & 0.58 & 0.76 & 1.16 \\
    1984-03-17 & 0.58 & 0.58 & 0.58 & 1.16 \\
    1979-09-16 & 2.09 & 2.09 & 2.09 & 2.09 \\
    1900-01-01 & 2.09 & 2.09 & 2.09 & 2.09 \\
    \bottomrule
  \end{tabular}
  }
\end{defaulttable}

\begin{defaulttable}
  \caption{\texttt{unknown} 최상층 지붕에 적용되는 열관류율
    $[\mathrm{W/(m^2K)}]$}
  \label{tbl:unknown-db-roof}
  {\small
  \renewcommand{\arraystretch}{1.05}
  \begin{tabular}{lrrrr}
    \toprule
    시행일자 & 중부1 & 중부2 & 남부 & 제주도 \\ \midrule
    2018-09-01 & 0.15 & 0.15 & 0.18 & 0.25 \\
    2016-01-01 & 0.15 & 0.15 & 0.18 & 0.25 \\
    2013-10-01 & 0.18 & 0.18 & 0.22 & 0.28 \\
    2013-09-01 & 0.18 & 0.18 & 0.22 & 0.28 \\
    2010-11-05 & 0.20 & 0.20 & 0.24 & 0.29 \\
    2008-07-10 & 0.29 & 0.29 & 0.35 & 0.41 \\
    2002-08-31 & 0.29 & 0.29 & 0.35 & 0.41 \\
    2001-01-17 & 0.29 & 0.29 & 0.35 & 0.41 \\
    1996-02-09 & 0.41 & 0.41 & 0.52 & 0.76 \\
    1992-06-01 & 0.41 & 0.41 & 0.52 & 0.76 \\
    1987-07-21 & 0.41 & 0.41 & 0.52 & 0.76 \\
    1984-03-17 & 0.58 & 0.58 & 0.58 & 1.16 \\
    1979-09-16 & 1.04 & 1.04 & 1.04 & 1.04 \\
    1900-01-01 & 1.04 & 1.04 & 1.04 & 1.04 \\
    \bottomrule
  \end{tabular}
  }
\end{defaulttable}

\begin{defaulttable}
  \caption{\texttt{unknown} 최하층 바닥에 적용되는 열관류율
    $[\mathrm{W/(m^2K)}]$}
  \label{tbl:unknown-db-groundfloor}
  {\small
  \renewcommand{\arraystretch}{1.05}
  \begin{tabular}{lrrrr}
    \toprule
    시행일자 & 중부1 & 중부2 & 남부 & 제주도 \\ \midrule
    2018-09-01 & 0.17 & 0.20 & 0.25 & 0.33 \\
    2016-01-01 & 0.22 & 0.22 & 0.25 & 0.33 \\
    2013-10-01 & 0.29 & 0.29 & 0.33 & 0.39 \\
    2013-09-01 & 0.29 & 0.29 & 0.29 & 0.29 \\
    2010-11-05 & 0.41 & 0.41 & 0.41 & 0.41 \\
    2008-07-10 & 0.41 & 0.41 & 0.47 & 0.52 \\
    2002-08-31 & 0.41 & 0.41 & 0.47 & 0.52 \\
    2001-01-17 & 0.41 & 0.41 & 0.47 & 0.52 \\
    1996-02-09 & 0.58 & 0.58 & 0.76 & 1.16 \\
    1992-06-01 & 0.58 & 0.58 & 0.76 & 1.16 \\
    1987-07-21 & 0.58 & 0.58 & 0.76 & 1.16 \\
    1984-03-17 & 0.58 & 0.58 & 0.58 & 1.16 \\
    1979-09-16 & 1.74 & 1.74 & 1.74 & 1.74 \\
    1900-01-01 & 1.74 & 1.74 & 1.74 & 1.74 \\
    \bottomrule
  \end{tabular}
  }
\end{defaulttable}

\begin{defaulttable}
  \caption{\texttt{unknown} 층간바닥에 적용되는 열관류율
    $[\mathrm{W/(m^2K)}]$}
  \label{tbl:unknown-db-interfloor}
  {\small
  \renewcommand{\arraystretch}{1.05}
  \begin{tabular}{lrrrr}
    \toprule
    시행일자 & 중부1 & 중부2 & 남부 & 제주도 \\ \midrule
    2018-09-01 & 1.16 & 1.16 & 1.16 & 1.16 \\
    2016-01-01 & 1.16 & 1.16 & 1.16 & 1.16 \\
    2013-10-01 & 1.16 & 1.16 & 1.16 & 1.16 \\
    2013-09-01 & 1.16 & 1.16 & 1.16 & 1.16 \\
    2010-11-05 & 1.16 & 1.16 & 1.16 & 1.16 \\
    2008-07-10 & 1.16 & 1.16 & 1.16 & 1.16 \\
    2002-08-31 & 1.16 & 1.16 & 1.16 & 1.16 \\
    2001-01-17 & 1.16 & 1.16 & 1.16 & 1.16 \\
    1996-02-09 & 0.58 & 0.58 & 0.76 & 1.16 \\
    1992-06-01 & 0.58 & 0.58 & 0.76 & 1.16 \\
    1987-07-21 & 0.58 & 0.58 & 0.76 & 1.16 \\
    1984-03-17 & 0.58 & 0.58 & 0.58 & 1.16 \\
    1979-09-16 & 1.74 & 1.74 & 1.74 & 1.74 \\
    1900-01-01 & 1.74 & 1.74 & 1.74 & 1.74 \\
    \bottomrule
  \end{tabular}
  }
\end{defaulttable}

\subsubsection{내장 기본 재료 DB}

기본 구조체 생성에는 \cref{tbl:default-materials}의 네 개 재료를 사용한다.

\begin{defaulttable}
  \caption{\simulator 의 내장 기본 재료 DB}
  \label{tbl:default-materials}
  {\small
  \begin{tabular}{lllrrr}
    \toprule
    DB 키 & 재료명 & 사용 위치
      & 열전도율
      & 밀도
      & 비열 \\
    & & &
      $[\mathrm{W/(mK)}]$
      & $[\mathrm{kg/m^3}]$
      & $[\mathrm{J/(kgK)}]$ \\ \midrule
    \texttt{concrete}
      & 콘크리트
      & 법규 구조체의 콘크리트층
      & 2.500 & 2400 & 880 \\

    \texttt{insulation}
      & 단열재
      & 법규 구조체의 외단열층
      & 0.050 & 30 & 1500 \\

    \texttt{gypsumboard}
      & 석고보드
      & 기본 내벽의 양면 마감층
      & 0.180 & 800 & 1090 \\

    \texttt{glasswool}
      & 글라스울
      & 기본 내벽의 중심 단열층
      & 0.035 & 19 & 960 \\
    \bottomrule
  \end{tabular}
  }
\end{defaulttable}

법규 DB에 의해 생성되는 구조체는 \texttt{insulation}과 \texttt{concrete}를 사용한다. 기본 내벽에는 \texttt{gypsumboard}와 \texttt{glasswool}을 사용한다.

\subsubsection{목표 열관류율의 재료층 변환}

선택된 목표 열관류율 $U_t$는 단열재와 콘크리트로 구성된 구조체로 변환한다. 실내·실외 표면열저항은 각각 다음 값을 사용한다.

\[
  R_{si}=0.110\ \mathrm{m^2K/W},
  \qquad
  R_{se}=0.043\ \mathrm{m^2K/W}
\]

콘크리트 두께 $t_c$의 초기값은 0.19 m이고, 단열재 두께는 다음 식으로 계산한다.

\begin{equation}
  t_{\mathrm{ins}}
  =
  k_{\mathrm{ins}}
  \left(
    \frac{1}{U_t}
    -
    R_{si}
    -
    R_{se}
    -
    \frac{t_c}{k_c}
  \right),
  \label{eq:unknown-insulation-thickness}
\end{equation}

여기서 내장 재료 DB에 따라

\[
  k_{\mathrm{ins}}=0.050\ \mathrm{W/(mK)},
  \qquad
  k_c=2.500\ \mathrm{W/(mK)}
\]

를 사용한다.

현재 \texttt{unknown} 해소 경로에서 열관류율은 0.15에서 2.09 $\mathrm{W/(m^2K)}$까지이다. 이 범위에서는 계산된 단열재 두께가 모두 양수이므로 콘크리트 두께는 항상 초기값인 190.0 mm로 유지된다. 현재 DB에서 사용되는 서로 다른 모든 $U_t$와 이에 대응하는 단열재 두께는 \cref{tbl:unknown-generated-insulation-thickness}와 같다. 표의 두께는 표시를 위해 0.1 mm로 반올림한 값이며, 실제 계산에는 반올림하지 않은 부동소수점 값을 사용한다. 이 때, 모든 법규 DB 기반 \texttt{unknown} 구조체는 외단열 구조체로 생성된다. 

\begin{defaulttable}
  \caption{내장 DB의 목표 열관류율별 생성 단열재 두께}
  \label{tbl:unknown-generated-insulation-thickness}
  {\small
  \begin{tabular}{rr@{\qquad}rr@{\qquad}rr}
    \toprule
    $U_t$ & $t_{\mathrm{ins}}$
      & $U_t$ & $t_{\mathrm{ins}}$
      & $U_t$ & $t_{\mathrm{ins}}$ \\
    $[\mathrm{W/(m^2K)}]$ & $[\mathrm{mm}]$
      & $[\mathrm{W/(m^2K)}]$ & $[\mathrm{mm}]$
      & $[\mathrm{W/(m^2K)}]$ & $[\mathrm{mm}]$ \\ \midrule
    0.15 & 321.9 & 0.29 & 161.0 & 0.45 & 99.7 \\
    0.17 & 282.7 & 0.32 & 144.8 & 0.47 & 94.9 \\
    0.18 & 266.3 & 0.33 & 140.1 & 0.52 & 84.7 \\
    0.20 & 238.6 & 0.34 & 135.6 & 0.58 & 74.8 \\
    0.22 & 215.8 & 0.35 & 131.4 & 0.76 & 54.3 \\
    0.24 & 196.9 & 0.36 & 127.4 & 1.04 & 36.6 \\
    0.25 & 188.5 & 0.39 & 116.8 & 1.16 & 31.7 \\
    0.26 & 180.9 & 0.41 & 110.5 & 1.74 & 17.3 \\
    0.27 & 173.7 & 0.43 & 104.8 & 2.09 & 12.5 \\
    0.28 & 167.1 & 0.44 & 102.2 &      &      \\
    \bottomrule
  \end{tabular}
  }
\end{defaulttable}

\subsection{\texttt{open} 구조체}

\inlinecode{construction_id}가 \inlinecode{open}이면 EnergyPlus의 \texttt{Construction:AirBoundary}를 생성한다. 공기교환 방식은 \texttt{SimpleMixing}, 존 간 혼합률은 0.5 ACH이다. 두 존은 별도 존으로 유지되고 공기경계 객체가 혼합을 계산한다.

\subsection{쿨루프 적용}

외기에 노출된 천장에 \inlinecode{coolroof_reflectance}가 입력되면 원 구조체의 가장 바깥쪽 재료층을 복제하고, 그 층의 태양 흡수율과 장파방사 흡수율을 모두 $1-\rho_{roof}$로 설정한다. 여기서 $\rho_{roof}$는 입력 반사율이다. 열전도율, 밀도, 비열, 두께 및 나머지 재료층은 변하지 않는다.

\subsection{창호와 문의 객체 변환}

투명 개구부 구조체는 입력 열관류율과 태양열취득계수를 사용하는 \texttt{WindowMaterial:}\allowbreak\texttt{Simple}\allowbreak\texttt{Glazing}\allowbreak\texttt{System}으로 변환한다.

불투명 문은 입력 열관류율의 역수를 열저항으로 사용하는 \texttt{Material:NoMass}로 생성한다.


% ---------------------------------------------------------------------------- %
%                                  NEW SECTION                                 %
% ---------------------------------------------------------------------------- %

\section{프로필}
\label{section:profile-conversion}

\simulator 의 용도프로필 데이터베이스는 용도별 대표값을 한 행에 저장한 정적 CSV로 구성된다.
존의 \inlinecode{profile} 값은 CSV의 \inlinecode{Name}과 일치하는 행을 선택한다.
선택된 행의 시간, 일일 발열량, 설정온도, 요일 플래그 및 휴가기간을 이용하여 재실, 냉난방 가용, 조명, 내부발열 및 급탕에 관한 연간 시계열을 생성한다.
기본 데이터는 국내 용도프로필 항목을 따른다 \cite{law:usageprofile}.

\subsection{정적 프로필 데이터베이스}

\begin{defaulttable}
  \caption{정적 용도프로필 CSV의 입력 항목}
  \label{tab:static-profile-fields}
  \small
  \begin{tabularx}{\textwidth}{p{3.0cm}p{4.8cm}X}
    \toprule
    구분 & CSV 열 & 저장 내용 \\ \midrule

    프로필 식별자 & \texttt{Name} & 존의 \inlinecode{profile} 값으로 조회하는 용도명 \\
    재실시간 & \texttt{Occupant-Start}, \texttt{Occupant-End} & 일간 재실 시작·종료시각 \\
    냉난방 가용시간 & \texttt{HVAC-Start}, \texttt{HVAC-End} & 존 공급설비의 일간 가용 시작·종료시각 \\
    프로필 외기량 & \texttt{Ventilation} & 용도별 면적당 외기량 속성 $[\mathrm{m^3/(m^2h)}]$ \\
    일일 급탕열량 & \texttt{DomesticHotWater} & 면적당 일일 급탕열량 $[\mathrm{Wh/(m^2d)}]$ \\
    일일 조명시간 & \texttt{LightingHours} & 일간 조명 가동분율의 적분시간 $[\mathrm{h}]$ \\
    일일 인체발열량 & \texttt{Occupancy} & 면적당 일일 인체 현열량 $[\mathrm{Wh/(m^2d)}]$ \\
    일일 기기발열량 & \texttt{Equipment} & 면적당 일일 기기발열량 $[\mathrm{Wh/(m^2d)}]$ \\
    설정온도 & \texttt{Heating-Setpoint}, \texttt{Cooling-Setpoint} & 난방 및 냉방 설정온도 $[{}^\circ\mathrm{C}]$ \\
    운전일 & \texttt{Monday}--\texttt{Sunday}, \texttt{Holiday} & 요일 및 공휴일별 운전 여부; 1은 운전, 0은 비운전 \\
    휴가기간 & \texttt{Vacations} & \texttt{MM/DD-MM/DD} 형식의 비운전 날짜 구간 \\
    \bottomrule
  \end{tabularx}
\end{defaulttable}

\subsection{내장 정적 프로필 DB의 전체 값}

\subsubsection{재실시간, 냉난방 가용시간 및 외기량}

현재 기본 및 확장 DB의 24개 프로필은 모두 난방 설정온도
20~$^{\circ}$C와 냉방 설정온도 26~$^{\circ}$C를 사용한다.
프로필별 시간 및 \texttt{Ventilation} 값은
\cref{tab:profile-db-time}와 같다.

표의 시간구간은 시작시각을 포함하고 종료시각을 포함하지 않는다.
\texttt{객실}의 21:00--08:00은 자정을 통과하는 시간구간이다.

\begin{defaulttable}
  \caption{내장 프로필 DB의 시간 및 외기량 값}
  \label{tab:profile-db-time}
  {\small
  \renewcommand{\arraystretch}{1.08}
  \begin{tabularx}{\textwidth}{Xccr}
    \toprule
    용도 & 재실시간 & 냉난방 가용시간 & \makecell[r]{\texttt{Ventilation}\\ $[\mathrm{m^3/(m^2h)}]$} \\
    \midrule
    주거공간             & 00:00--24:00 & 00:00--24:00 & 1.1 \\
    소규모사무실         & 09:00--18:00 & 07:00--18:00 & 4 \\
    대규모사무실         & 09:00--18:00 & 07:00--18:00 & 6 \\
    회의실 및 세미나실   & 07:00--18:00 & 07:00--18:00 & 15 \\
    강당                 & 07:00--18:00 & 07:00--18:00 & 2 \\
    구내식당             & 08:00--15:00 & 08:00--15:00 & 18 \\
    화장실               & 07:00--18:00 & 07:00--18:00 & 15 \\
    그 외 체류공간       & 07:00--18:00 & 07:00--18:00 & 7 \\
    부속공간             & 07:00--18:00 & 07:00--18:00 & 0.15 \\
    창고/설비/문서실     & 07:00--18:00 & 07:00--18:00 & 0.15 \\
    전산실               & 00:00--24:00 & 00:00--24:00 & 1.3 \\
    주방 및 조리실       & 08:00--15:00 & 08:00--15:00 & 90 \\
    병실                 & 00:00--24:00 & 00:00--24:00 & 4 \\
    객실                 & 21:00--08:00 & 21:00--08:00 & 3 \\
    교실(초중고)         & 08:00--15:00 & 08:00--15:00 & 10 \\
    강의실(대학)         & 09:00--18:00 & 09:00--18:00 & 30 \\
    매장(상점/백화점)    & 08:00--20:00 & 08:00--20:00 & 4 \\
    전시실(전시관/박물관) & 10:00--18:00 & 10:00--18:00 & 2 \\
    열람실(도서관)       & 08:00--20:00 & 08:00--20:00 & 8 \\
    체육시설             & 08:00--23:00 & 08:00--23:00 & 3 \\
    구내식당(초중고)     & 11:00--15:00 & 11:00--15:00 & 18 \\
    주방 및 조리실(초중고) & 08:00--15:00 & 08:00--15:00 & 90 \\
    체육시설(초중고)     & 09:00--16:00 & 09:00--16:00 & 3 \\
    교실(어린이집)       & 08:00--15:00 & 08:00--15:00 & 10 \\

    \bottomrule
  \end{tabularx}
  }
\end{defaulttable}

\texttt{Ventilation} 값은
\texttt{KoreanUsageProfile} 객체의 속성으로 저장된다. 현재 프로필 변환에서 생성되는 재실자, 기기, 조명, 급탕, 냉난방 가용 및 설정온도 스케줄에는 이 값이 사용되지 않는다. 존의 EnergyPlus 환기 객체는 존의 환기설비 연결정보에 따라 별도로 생성한다.

\subsubsection{일일 급탕, 조명, 인체발열 및 기기발열}

\cref{tab:profile-db-load}는 csv에 저장된 면적당 일일 발열량을 나타낸다. 이 값들은 그 자체가 EnergyPlus의 시간별 부하밀도는 아니며, 재실시간으로 나누어 시간별 값으로 변환된다.

\begin{defaulttable}
  \caption{내장 프로필 DB의 일일 발열량 값}
  \label{tab:profile-db-load}
  {\small
  \renewcommand{\arraystretch}{1.08}
  \begin{tabularx}{\textwidth}{Xrrrr}
    \toprule
    용도 & \makecell[r]{급탕열량 \\ $[\mathrm{Wh/(m^2d)}]$} & \makecell[r]{조명시간\\$[\mathrm{h}]$} & \makecell[r]{인체발열량\\ $[\mathrm{Wh/(m^2d)}]$} & \makecell[r]{기기발열량\\ $[\mathrm{Wh/(m^2d)}]$} \\ 
    \midrule

    주거공간                & 84   & 5  & 53   & 52   \\
    소규모사무실            & 30   & 6  & 30   & 42   \\
    대규모사무실            & 30   & 9  & 55.8 & 126  \\
    회의실 및 세미나실      & 30   & 11 & 96   & 8    \\
    강당                    & 30   & 11 & 36   & 24   \\
    구내식당                & 1250 & 7  & 177  & 10   \\
    화장실                  & 0    & 11 & 0    & 0    \\
    그 외 체류공간          & 30   & 11 & 96   & 8    \\
    부속공간                & 0    & 11 & 0    & 0    \\
    창고/설비/문서실        & 0    & 11 & 0    & 0    \\
    전산실                  & 30   & 12 & 15   & 1800 \\
    주방 및 조리실          & 0    & 7  & 56   & 1800 \\
    병실                    & 82   & 12 & 108  & 24   \\
    객실                    & 82   & 4  & 70   & 44   \\
    교실(초중고)            & 30   & 6  & 100  & 20   \\
    강의실(대학)            & 30   & 6  & 420  & 24   \\
    매장(상점/백화점)       & 30   & 12 & 84   & 24   \\
    전시실(전시관/박물관)   & 30   & 8  & 28   & 0    \\
    열람실(도서관)          & 30   & 12 & 168  & 0    \\
    체육시설                & 220  & 15 & 60   & 0    \\
    구내식당(초중고)        & 1250 & 4  & 177  & 10   \\
    주방 및 조리실(초중고)  & 0    & 7  & 56   & 1800 \\
    체육시설(초중고)        & 220  & 7  & 60   & 0    \\
    교실(어린이집)          & 30   & 6  & 100  & 20   \\

    \bottomrule
  \end{tabularx}
  }
\end{defaulttable}

\subsubsection{요일 및 공휴일 운전 플래그}

현재 DB에는 세 종류의 운전일 조합이 있다.
\cref{tab:profile-db-operating-days}의 운전일은 해당 CSV 플래그가 1인 날짜를 나타내며, 표에 없는 요일과 공휴일 플래그는 0이다.

\begin{defaulttable}
  \caption{내장 프로필 DB의 요일 및 공휴일 운전 플래그}
  \label{tab:profile-db-operating-days}
  {\small
  \renewcommand{\arraystretch}{1.12}
  \begin{tabularx}{\textwidth}{p{2.5cm}p{4.5cm}X}
    \toprule
    운전 유형 & 값이 1인 플래그 & 해당 프로필 \\ \midrule

    연중 운전
      & 월,화,수,목,금,토,일,공휴일
      & 주거공간, 전산실, 병실, 객실 \\

    평일 운전
      & 월,화,수,목,금
      & 소규모사무실, 대규모사무실, 회의실 및 세미나실,
        강당, 구내식당, 화장실, 그 외 체류공간, 부속공간,
        창고/설비/문서실, 주방 및 조리실, 교실(초중고),
        강의실(대학), 전시실(전시관/박물관),
        구내식당(초중고), 주방 및 조리실(초중고),
        교실(어린이집) \\

    월-토 운전
      & 월,화,수,목,금,토
      & 매장(상점/백화점), 열람실(도서관), 체육시설,
        체육시설(초중고) \\

    \bottomrule
  \end{tabularx}
  }
\end{defaulttable}

\subsubsection{휴가기간}

현재 DB에서 \texttt{Vacations} 값이 존재하는 프로필은 네 개이다. 기간의 시작일과 종료일은 모두 휴가기간에 포함된다. 표에 포함되지 않은 나머지 20개 프로필의 \texttt{Vacations} 값은 비어 있다.

\begin{defaulttable}
  \caption{내장 프로필 DB의 휴가기간}
  \label{tab:profile-db-vacations}
  {\small
  \begin{tabularx}{\textwidth}{p{5.0cm}X}
    \toprule
    프로필 & 휴가기간 \\ \midrule

    교실(초중고)
      & \texttt{01/01-02/14};
        \texttt{07/22-08/27};
        \texttt{12/22-12/31} \\

    강의실(대학)
      & \texttt{01/01-02/28};
        \texttt{07/08-08/31};
        \texttt{12/14-12/31} \\

    구내식당(초중고)
      & \texttt{01/01-02/14};
        \texttt{07/22-08/27};
        \texttt{12/22-12/31} \\

    주방 및 조리실(초중고)
      & \texttt{01/01-02/14};
        \texttt{07/22-08/27};
        \texttt{12/22-12/31} \\

    \bottomrule
  \end{tabularx}
  }
\end{defaulttable}

\subsection{시간별 변환 결과}

\cref{tab:profile-derived-value}는
현재 DB의 정적 일일 발열량을 변환한 운전시간 중 일정값이다.

표시값은 소수점 다섯째 자리까지 반올림하였다. 실제 스케줄 생성에는
반올림하지 않은 부동소수점 값을 사용한다.

\begin{defaulttable}
  \caption{현재 내장 프로필의 시간별 변환 결과}
  \label{tab:profile-derived-value}
  {\small
  \renewcommand{\arraystretch}{1.08}
  \begin{tabularx}{\textwidth}{Xrrrr}
    \toprule
    용도
      & $H_{occ}$
      & \makecell[r]{$n_{occ}$\\$[\mathrm{person/m^2}]$}
      & \makecell[r]{$p_{equip}$\\$[\mathrm{W/m^2}]$}
      & \makecell[r]{$v_{dhw}$\\$[\mathrm{L/(m^2h)}]$} \\ \midrule

    주거공간                & 24 & 0.03155 & 2.16667   & 0.08750 \\
    소규모사무실            & 9  & 0.04762 & 4.66667   & 0.08333 \\
    대규모사무실            & 9  & 0.08857 & 14.00000  & 0.08333 \\
    회의실 및 세미나실      & 11 & 0.12468 & 0.72727   & 0.06818 \\
    강당                    & 11 & 0.04675 & 2.18182   & 0.06818 \\
    구내식당                & 7  & 0.36122 & 1.42857   & 4.46429 \\
    화장실                  & 11 & 0       & 0         & 0       \\
    그 외 체류공간          & 11 & 0.12468 & 0.72727   & 0.06818 \\
    부속공간                & 11 & 0       & 0         & 0       \\
    창고/설비/문서실        & 11 & 0       & 0         & 0       \\
    전산실                  & 24 & 0.00893 & 75.00000  & 0.03125 \\
    주방 및 조리실          & 7  & 0.11429 & 257.14286 & 0       \\
    병실                    & 24 & 0.06429 & 1.00000   & 0.08542 \\
    객실                    & 11 & 0.09091 & 4.00000   & 0.18636 \\
    교실(초중고)            & 7  & 0.20408 & 2.85714   & 0.10714 \\
    강의실(대학)            & 9  & 0.66667 & 2.66667   & 0.08333 \\
    매장(상점/백화점)       & 12 & 0.10000 & 2.00000   & 0.06250 \\
    전시실(전시관/박물관)   & 8  & 0.05000 & 0         & 0.09375 \\
    열람실(도서관)          & 12 & 0.20000 & 0         & 0.06250 \\
    체육시설                & 15 & 0.05714 & 0         & 0.36667 \\
    구내식당(초중고)        & 4  & 0.63214 & 2.50000   & 7.81250 \\
    주방 및 조리실(초중고)  & 7  & 0.11429 & 257.14286 & 0       \\
    체육시설(초중고)        & 7  & 0.12245 & 0         & 0.78571 \\
    교실(어린이집)          & 7  & 0.20408 & 2.85714   & 0.10714 \\

    \bottomrule
  \end{tabularx}
  }
\end{defaulttable}

예를 들어 소규모사무실의 일일 인체발열량 30~Wh/(m$^2$d)는
9시간의 재실시간과 70~W/person을 이용하여
0.04762~person/m$^2$로 변환된다. 기기발열량 42~Wh/(m$^2$d)는
4.66667~W/m$^2$, 급탕열량 30~Wh/(m$^2$d)는
0.08333~L/(m$^2$h)로 변환된다.

\begin{defaultfigure}
  \includegraphics[width=\textwidth]{profile_small_office_daily.png}
  \caption{소규모사무실 운전일의 재실·냉난방 가용·조명 스케줄}
  \label{fig:profile-small-office-daily}
\end{defaultfigure}

\subsection{조명시간의 10분 구간 배치}
\label{subsection:profile-lighting-allocation}

CSV에는 조명 시작시각과 종료시각이 아니라 일일 조명시간
$H_{light}$가 저장되어 있다. 조명 스케줄 $f_{light,i}$는
0--1의 가동분율이며, 운전일의 적분값이
\texttt{LightingHours}와 같도록 생성한다.

\begin{equation}
  \sum_{i=1}^{144}
  f_{light,i}
  \left(\frac{1}{6}\right)
  =
  H_{light}.
  \label{eq:profile-lighting-integral}
\end{equation}

조명시간의 배치 절차는 다음과 같다.

\begin{enumerate}
  \item 재실시간에 포함되는 10분 구간을 후보로 만든다.
  \item 각 구간의 중앙시각과 12:00 사이의 순환시간거리를 계산한다.
  \item 00:00--06:00에 속하지 않는 후보를 시간거리가 큰 순서로
        선택한다.
  \item 같은 시간거리를 갖는 구간에 남은 조명시간이 걸리면
        동일한 가동분율을 배분한다.
  \item 조명시간이 남으면 00:00--06:00의 재실 구간을
        00:00부터 순서대로 선택한다.
\end{enumerate}

현재 DB에서 생성되는 일간 조명시간대는
\cref{tab:profile-lighting-db-results}와 같다.

\begin{defaulttable}
  \caption{현재 내장 프로필에서 생성되는 일간 조명시간대}
  \label{tab:profile-lighting-db-results}
  {\small
  \renewcommand{\arraystretch}{1.12}
  \begin{tabularx}{\textwidth}{p{4.2cm}rX}
    \toprule
    조명시간대
      & $H_{light}$ [h]
      & 해당 프로필 \\ \midrule

    19:00--24:00
      & 5
      & 주거공간 \\

    09:00--10:30, 13:30--18:00
      & 6
      & 소규모사무실, 강의실(대학) \\

    09:00--18:00
      & 9
      & 대규모사무실 \\

    07:00--18:00
      & 11
      & 회의실 및 세미나실, 강당, 화장실,
        그 외 체류공간, 부속공간, 창고/설비/문서실 \\

    08:00--15:00
      & 7
      & 구내식당, 주방 및 조리실,
        주방 및 조리실(초중고) \\

    06:00--09:00, 15:00--24:00
      & 12
      & 전산실, 병실 \\

    06:00--07:00, 21:00--24:00
      & 4
      & 객실 \\

    08:00--11:30, 12:30--15:00
      & 6
      & 교실(초중고), 교실(어린이집) \\

    08:00--20:00
      & 12
      & 매장(상점/백화점), 열람실(도서관) \\

    10:00--18:00
      & 8
      & 전시실(전시관/박물관) \\

    08:00--23:00
      & 15
      & 체육시설 \\

    11:00--15:00
      & 4
      & 구내식당(초중고) \\

    09:00--16:00
      & 7
      & 체육시설(초중고) \\

    \bottomrule
  \end{tabularx}
  }
\end{defaulttable}

\begin{defaultfigure}
  \includegraphics[width=\textwidth]{profile_lighting_examples.png}
  \caption{용도별 일일 조명시간의 10분 구간 배치 예}
  \label{fig:profile-lighting-examples}
\end{defaultfigure}

\subsection{휴가기간의 연간 적용}

휴가기간은 2026년의 월·일 조합으로 변환되며, 각 구간의 시작일과
종료일을 모두 포함한다. 휴가기간에는 해당 날짜의 재실, 기기발열,
급탕, 냉난방 가용 및 조명 스케줄이 0이 된다.

\begin{defaultfigure}
  \includegraphics[width=\textwidth]{profile_classroom_vacation_mask.png}
  \caption{교실(초중고) 프로필의 2026년 휴가기간 마스크}
  \label{fig:profile-classroom-vacation}
\end{defaultfigure}

\subsection{EnergyPlus 객체 연결}

생성된 연간 스케줄은 EnergyPlus의 \texttt{Schedule:Compact} 객체로 변환한다. 기간은 \texttt{Through}, 요일 규칙은 \texttt{For}, 일간 값은 \texttt{Until} 필드로 기록한다.

\begin{defaulttable}
  \caption{동적 프로필 스케줄의 EnergyPlus 객체 연결}
  \label{tab:profile-energyplus-connection}
  \small
  \begin{tabularx}{\textwidth}{p{3.4cm}p{4.4cm}X}
    \toprule
    스케줄 & 참조 객체 & 적용 필드 또는 방식 \\ \midrule

    난방·냉방 설정온도
      & \makecell[l]{\texttt{ThermostatSetpoint:}\\
                     \texttt{DualSetpoint}}
      & 난방 및 냉방 설정온도 스케줄명 \\

    냉난방 가용
      & 존 공급설비
      & 장비 가용스케줄 \\

    재실자
      & \texttt{People}
      & 원 스케줄을 최댓값으로 정규화하고,
        최댓값을 \texttt{People/Area}에 입력 \\

    기기발열
      & \texttt{ElectricEquipment}
      & 원 스케줄을 최댓값으로 정규화하고,
        최댓값을 \texttt{Watts/Area}에 입력 \\

    조명
      & \texttt{Lights}
      & 가동분율 스케줄과 존의
        \inlinecode{light_density}를 사용 \\

    급탕
      & \texttt{Schedule:Compact}
      & 프로필의 급탕 연간 스케줄로 출력 \\

    \bottomrule
  \end{tabularx}
\end{defaulttable}


% ---------------------------------------------------------------------------- %
%                                  NEW SECTION                                 %
% ---------------------------------------------------------------------------- %

\section{냉난방설비}
\label{section:hvac-conversion}

\simulator 의 냉난방설비 입력은  열을 생산하는 \textbf{생산설비}(source system), 생산된 열 또는 냉열을 존에 전달하는 \textbf{공급설비}(supply system), 그리고 설비를 사용하는 \textbf{존}(zone)의 세 계층을 식별자로 연결하는 방식이다. 입력 항목과 허용값은 \cref{subsubsection:ioref:supplysystem,subsubsection:ioref:sourcesystem}에 정의되어 있다. 변환 과정에서는 이 관계를 먼저 \pymodule 객체 그래프로 해석한 뒤, \texttt{idragon} 모듈이 EnergyPlus 열원, PlantLoop, 존 말단기기, 노드 및 연결 객체로 전개한다. 각 EnergyPlus 객체와 연결 필드의 의미는 공식 입력출력 참조를 기준으로 한다 \cite{energyplusior242}.

\subsection{변환 계층과 객체 생성 순서}

입력의 \inlinecode{name}은 사용자 표시용 속성이고, 객체 간 참조와 중복 제거에는 \inlinecode{id}가 사용된다. \pymodule은 생산설비와 공급설비를 각각 ID 기반 사전으로 생성한다. 공급설비의 \texttt{source\_}\allowbreak\texttt{system\_}\allowbreak\texttt{id}는 실제 생산설비 객체로 치환한다. 각 존의 난방 및 냉방 공급설비 ID도 같은 방식으로 공급설비 객체에 연결한다. 이후 \pymodule 객체의 ID가 \texttt{idragon} 객체명으로 전달되므로, 최종 IDF의 장비명과 노드명은 대체로 입력 ID에서 파생된다.

구체적인 생성 순서는 다음과 같다.

\begin{enumerate}
  \item 입력의 모든 생산설비를 먼저 읽어 \inlinecode{id}를 키로 하는 생산설비 사전을 만든다.
  \item 공급설비를 생성하면서 \inlinecode{source_system_id}를 생산설비 객체로 치환한다. 독립형 전기설비와 패키지형 에어컨은 생산설비를 요구하지 않는다.
  \item 존을 생성하면서 난방용, 냉방용 공급설비 ID를 실제 공급설비 객체로 치환한다. 이때 공급설비 종류와 생산설비 종류가 해당 운전모드를 지원하는지 검사한다.
  \item 동일한 생산설비 ID가 여러 존에서 참조되더라도 생산설비 및 주 PlantLoop는 한 번만 만든다.
  \item 각 존에 크기산정 객체(\texttt{Sizing:Zone}), 존 장비목록, 장비연결, 온도조절기 및 이중 설정온도 객체를 만든다.
  \item 생산설비 IDF 객체와 주 PlantLoop를 만든 뒤, 존별 공급설비 객체를 생성한다. 한 존에서 동일한 공급설비가 난방과 냉방에 모두 지정되면 하나의 말단기기를 양 모드로 생성하고, 서로 다른 공급설비가 지정되면 두 말단기기를 별도로 생성한다.
  \item 공급설비가 반환한 후처리기를 실행하여 존 공기 입·출구 노드, 존 장비목록, VRF 말단기 목록, 수배관 PlantLoop의 수요측 Branch 목록과 Splitter/Mixer를 갱신한다.
\end{enumerate}

공급설비가 하나도 지정되지 않은 존은 비공조 존으로 취급한다. 이 경우 계산 안정성을 위해 난방 및 냉방 설정온도가 각각 $-30\,^{\circ}\mathrm{C}$와 $50\,^{\circ}\mathrm{C}$인 \texttt{HVACTemplate:Zone:IdealLoadsAirSystem}을 생성하므로, 정상 기상조건에서는 실질적인 냉난방을 수행하지 않는다.

\subsection{공급설비와 생산설비의 연결 규칙}

\cref{tab:hvac-connectability}는 입력에서 허용되는 공급설비--생산설비 조합과 실제 EnergyPlus 말단기기 유형을 정리한 것이다. 한 개의 \inlinecode{fan_coil_unit} 입력은 한 개의 생산설비만 참조하므로, 동일 존에서 난방과 냉방을 모두 팬코일로 모사할 때에는 난방원과 냉열원을 각각 참조하는 두 개의 팬코일 입력을 생성하여 존의 난방 및 냉방 ID에 따로 지정해야 한다.

\begin{defaulttable}
  \caption{공급설비별 허용 생산설비와 EnergyPlus 구현 유형}
  \label{tab:hvac-connectability}
  \small
  \begin{tabularx}{\textwidth}{>{\raggedright\arraybackslash}p{4.5cm}>{\raggedright\arraybackslash}p{4.5cm}X}
    \toprule
    공급설비 입력 유형 & 허용 운전 및 생산설비 & EnergyPlus 구현 유형 \\ \midrule
    \makecell[l]{\texttt{packaged\_air\_}\\\texttt{conditioner}} & 냉방 / 생산설비 없음 & 전용 VRF 생산설비와 VRF 존 말단기 \\
    \makecell[l]{\texttt{air\_handling\_}\\\texttt{unit}} & 난방·냉방 / \texttt{heatpump}, \makecell[l]{\texttt{geothermal\_}\\\texttt{heatpump}} & VRF 존 말단기 \\
    \texttt{fan\_coil\_unit} & 난방 / \texttt{boiler}, \texttt{district\_heating} & 4관식 팬코일의 온수코일 사용 \\
    \texttt{fan\_coil\_unit} & 냉방 / \texttt{chiller}, \makecell[l]{\texttt{absorption\_}\\\texttt{chiller}} & 4관식 팬코일의 냉수코일 사용 \\
    \texttt{radiator} & 난방 / \texttt{boiler}, \texttt{district\_heating} & 온수식 방열기 \\
    \makecell[l]{\texttt{electric\_}\\\texttt{radiator}} & 난방 / 생산설비 없음 & 전기식 방열기 \\
    \texttt{radiant\_floor} & 난방 / \texttt{boiler}, \texttt{district\_heating} & 수식 저온 바닥복사난방 \\
    \makecell[l]{\texttt{electric\_radiant\_}\\\texttt{floor}} & 난방 / 생산설비 없음 & 전기식 저온 바닥복사난방 \\ \bottomrule
  \end{tabularx}
\end{defaulttable}

\subsection{생산설비의 EnergyPlus 변환}

\subsubsection{히트펌프 및 지열히트펌프}

\inlinecode{heatpump}는 \texttt{AirConditioner:VariableRefrigerantFlow} 객체로 변환된다. 냉방·난방 COP를 생략하면 각각 3.0을 사용한다. 입력의 COP와 용량은 각각 정격 COP와 정격 용량에 전달되며, 용량이 생략되면 \texttt{Autosize}가 적용된다. 온도 및 부분부하에 따른 용량비와 EIR 보정은 사용자가 별도 곡선을 입력하는 방식이 아니라, \texttt{idragon}에 내장된 저온·고온 및 부분부하 성능곡선 집합을 자동 생성하여 적용한다. 하나의 VRF 생산설비에는 \texttt{ZoneTerminalUnitList}가 생성되고, 이 생산설비를 참조하는 모든 존 말단기가 후처리 단계에서 목록에 추가된다.

\inlinecode{geothermal_heatpump}는 변환 시 일반 \inlinecode{heatpump}와 동일한 객체를 생성한다. 따라서 지중열교환기, 보어홀, 지중온도 응답 또는 수측 루프는 생성되지 않는다.

\subsubsection{전기식 냉동기}

\inlinecode{chiller}는 \texttt{Chiller:Electric:EIR}와 냉수 PlantLoop 및 응축수 CondenserLoop의 세 부분으로 전개된다. COP를 생략하면 3.0, 용량을 생략하면 \texttt{Autosize}를 사용한다. 압축기 유형 \inlinecode{turbo}, \inlinecode{screw}, \inlinecode{reciprocating}은 각각 다른 내장 성능곡선 집합을 선택한다. 냉수 루프에는 가변속 펌프, 공급측 냉동기 Branch, 공급 및 수요 바이패스, Splitter/Mixer, 운전제어 및 설정온도 관리자가 생성된다. 냉수 공급 설정온도는 $6\,^{\circ}\mathrm{C}$, 설계 온도차는 $4\,\mathrm{K}$이며 펌프 모터효율은 0.9로 고정된다.

응축수 루프에는 별도의 가변속 펌프와 냉각탑 또는 밀폐형 유체냉각기가 생성된다. 설정온도는 외기온도 추종 방식으로 계산하며, 외기온도 대비 $1.5\,\mathrm{K}$ 오프셋과 $20$--$50\,^{\circ}\mathrm{C}$ 제한을 적용한다. 응축수 루프의 설계 출구온도와 설계 온도차는 각각 $29\,^{\circ}\mathrm{C}$와 $5\,\mathrm{K}$이다. 냉각탑 입력 조합은 \cref{tab:coolingtower-mapping}과 같이 EnergyPlus 객체에 대응한다.

\begin{defaulttable}
  \caption{냉각탑 입력 조합의 EnergyPlus 객체 대응}
  \label{tab:coolingtower-mapping}
  \begin{tabular}{ccc}
    \toprule
    \inlinecode{coolingtower_type} & \inlinecode{coolingtower_control} & EnergyPlus 객체 \\ \midrule
    \texttt{open} & \texttt{single-speed} & \texttt{CoolingTower:SingleSpeed} \\
    \texttt{open} & \texttt{two-speed} & \texttt{CoolingTower:TwoSpeed} \\
    \texttt{closed} & \texttt{single-speed} & \texttt{FluidCooler:SingleSpeed} \\
    \texttt{closed} & \texttt{two-speed} & \texttt{FluidCooler:TwoSpeed} \\ \bottomrule
  \end{tabular}
\end{defaulttable}

냉각탑 용량이 생략되면 냉동기 용량을 사용하고, 냉동기 용량도 생략된 경우에는 냉각탑의 명목용량에 $1{,}000{,}000\,\mathrm{W}$를 임시값으로 부여한다. 실제 설계 유량, 팬동력, UA 등은 가능한 항목에 \texttt{Autosize} 또는 \texttt{Autocalculate}를 적용한다.

\subsubsection{흡수식 냉동기}

\inlinecode{absorption_chiller}는 \texttt{Chiller:Absorption}과 냉수 PlantLoop를 생성한다. 입력 COP의 기본값은 0.9이고, 용량을 생략하면 \texttt{Autosize}를 사용한다. 사용자가 보조 열원과 냉각탑을 입력하지 않아도, 천연가스 \texttt{Boiler:HotWater}와 개방형 냉각탑 \texttt{CoolingTower:SingleSpeed}을 자동 생성하고 재생기 측과 응축기 측에 연결한다. 보일러 효율의 기본값은 0.85이고, 냉각탑 용량은 흡수식 냉동기 용량을 따른다. 냉수 루프의 설정온도와 설계 온도차는 전기식 냉동기와 동일하게 $6\,^{\circ}\mathrm{C}$와 $4\,\mathrm{K}$이다.

\subsubsection{보일러 및 지역난방}

\inlinecode{boiler}는 온수보일러 객체(\texttt{Boiler:}\allowbreak\texttt{HotWater})와 온수 PlantLoop로 변환된다. 효율을 생략하면 0.85를 사용한다. 보일러 용량을 생략하면 \texttt{Autosize}를 사용하고, 입력 효율은 명목 열효율에 직접 전달한다. 온수 루프는 가변속 펌프, 공급측 보일러 Branch, 바이패스, Splitter/Mixer, 수요측 BranchList와 설정온도 관리자를 포함한다. 운전 설정온도는 $60\,^{\circ}\mathrm{C}$이며, \texttt{Sizing:Plant}에는 설계 출구온도 $80\,^{\circ}\mathrm{C}$와 설계 온도차 $10\,\mathrm{K}$가 적용된다.
ㄴ
\inlinecode{district_heating}은 별도 지역난방 객체가 아니라 온수보일러와 동일한 온수 루프로 등가화한다. 이때 효율은 1.0, 용량은 \texttt{Autosize}, 연료는 \texttt{OtherFuel1}이다.

\subsection{존 공급설비의 객체 전개}

\subsubsection{패키지형 에어컨과 공기조화기}

\inlinecode{packaged_air_conditioner}는 생산설비 ID를 요구하지 않는다. 냉방 COP를 생략하면 3.0을 사용하고, 냉방용량을 생략하면 \texttt{Autosize}를 적용한다. 변환 과정에서 해당 공급설비 전용의 전기식 VRF 외기측 장치를 내부적으로 만들고, 난방용량을 $10^{-3}\,\mathrm{W}$로 두어 난방을 사실상 비활성화한 뒤 입력 냉방 COP 및 냉방용량을 적용한다. 따라서 최종 IDF에는 패키지형 단일장치 객체가 아니라 \texttt{AirConditioner:}\allowbreak\texttt{Variable}\allowbreak\texttt{Refrigerant}\allowbreak\texttt{Flow}와 \texttt{ZoneHVAC:}\allowbreak\texttt{TerminalUnit:}\allowbreak\texttt{Variable}\allowbreak\texttt{Refrigerant}\allowbreak\texttt{Flow}의 조합이 생성된다.

\inlinecode{air_handling_unit}는 VRF 존 말단기로 변환된다. 각 말단기는 냉방 DX 코일, 난방 DX 코일 및 정풍량 팬을 포함한다. 냉방코일의 정격 현열비는 0.7, 팬 전효율은 0.7, 압력상승은 $100\,\mathrm{Pa}$, 모터효율은 0.9이다. 코일 정격풍량과 팬 최대풍량은 존 바닥면적 $A_z$에 대해 다음과 같이 정한다.

\begin{equation}
  \dot V_{\mathrm{VRF}} = 0.01 A_z \qquad [\mathrm{m^3/s}]
  \label{eq:vrf-terminal-airflow}
\end{equation}

말단기의 외기유량 필드는 냉방, 난방 및 무부하 상태 모두 0으로 두며, 환기는 \cref{section:ventilation-conversion}의 별도 존 환기 객체가 담당한다. 사용하지 않는 냉방 또는 난방 코일은 가용스케줄을 \texttt{ALLOFF}로 설정하고 정격용량을 $0.1\,\mathrm{W}$로 둔다. 팬은 Draw-Through 방식이며 운전·정지 전력은 각각 30 W와 20 W이다.

\subsubsection{팬코일유닛}

\inlinecode{fan_coil_unit}는 4관식 팬코일(\texttt{ZoneHVAC:FourPipeFanCoil})과 외기혼합기, 단속 팬, 냉수코일 및 온수코일로 전개된다. 제어방식은 \texttt{ConstantFanVariableFlow}이고, 팬 전효율 0.7, 압력상승 $100\,\mathrm{Pa}$, 모터효율 0.9를 사용한다.

EnergyPlus의 4관식 팬코일 객체는 냉수코일과 온수코일을 모두 요구하지만, \simulator 의 팬코일 입력은 하나의 생산설비만 참조한다. 이에 따라 난방용 팬코일에는 실제 온수 Branch와 함께 용량 및 효율이 $10^{-10}$인 비사용 냉동기·냉각탑 루프를 만들고, 냉방용 팬코일에는 실제 냉수 Branch와 함께 용량 및 효율이 $10^{-10}$인 비사용 보일러 루프를 만든다. 비사용 루프의 가용스케줄은 \texttt{ALLOFF}로 설정된다. 이 보조 루프는 EnergyPlus 객체 연결 요건을 만족시키기 위한 것으로 에너지 공급을 의도하지 않는다.

\subsubsection{방열기 및 바닥복사난방}

\inlinecode{radiator}는 온수 방열기와 설계 객체를 만들고, 해당 온수 Branch를 보일러 또는 지역난방 PlantLoop의 수요측에 추가한다. 난방용량을 생략하면 \texttt{Autosize}를 사용하고, 최대 온수유량도 \texttt{Autosize}한다. 복사분율은 기본값 0을 사용하므로 전열은 대류 성분으로 계산된다. \inlinecode{electric_radiator}는 동일한 형상의 전기식 방열기로 변환하며 효율 1.0, 복사분율 0을 사용하고 PlantLoop를 만들지 않는다.

\inlinecode{radiant_floor}는 존에 속한 모든 바닥면을 저온복사 표면그룹으로 묶는다. 이 표면그룹의 EnergyPlus 객체 유형은 \texttt{ZoneHVAC:LowTemperatureRadiant:SurfaceGroup}이다. 각 면의 유량분율은 해당 면적을 존 총 바닥면적으로 나눈 값으로 정한다. 각 바닥 구조체에는 내부 열원 속성을 만들고, 열원 위치를 구조체의 마지막 층 바로 앞에 배치하며 관 간격은 0.3 m로 둔다. 내부 열원 속성의 객체 유형은 \texttt{ConstructionProperty:InternalHeatSource}이다. 온수식 바닥복사난방은 가변유량 저온복사 객체와 온수 Branch를 생성하고, 전기식은 전기 저온복사 객체를 생성한다. 두 유형 모두 난방 설정온도 스케줄을 제어온도로 사용하고 제어 비례대역은 $2\,\mathrm{K}$, 제어방식은 \texttt{ZeroFlowPower}이다.

\subsection{공통 제어, 스케줄 및 자동용량산정}

공조 존에는 프로필에서 생성한 난방·냉방 설정온도 스케줄을 참조하는 이중 설정온도 객체인 \texttt{ThermostatSetpoint:}\allowbreak\texttt{Dual}\allowbreak\texttt{Setpoint}를 생성한다. 공급설비의 가용스케줄은 프로필의 냉난방 가동 스케줄을 사용한다. 용량 입력이 없는 열원, 코일, 방열기와 유량 입력이 없는 펌프·코일·팬은 가능한 항목에 \texttt{Autosize}를 사용한다.

존 설계조건은 입력 항목으로 노출하지 않고 \cref{tab:zone-sizing-assumptions}과 같이 고정한다. 이 값은 EnergyPlus 크기산정 단계에 사용되며, 실제 운전 설정온도는 프로필의 설정온도 스케줄에 따른다.

\begin{defaulttable}
  \caption{냉난방 존 자동용량산정의 고정 설계조건}
  \label{tab:zone-sizing-assumptions}
  \begin{tabular}{lcc}
    \toprule
    항목 & 난방 & 냉방 \\ \midrule
    설계 급기온도 & $50\,^{\circ}\mathrm{C}$ & $14\,^{\circ}\mathrm{C}$ \\
    설계 급기온도차 & $10\,\mathrm{K}$ & $10\,\mathrm{K}$ \\
    설계 급기 절대습도 & $0.004\,\mathrm{kg/kg}$ & $0.009\,\mathrm{kg/kg}$ \\
    크기산정 여유계수 & 1.25 & 1.15 \\ \bottomrule
  \end{tabular}
\end{defaulttable}

\section{환기설비}
\label{section:ventilation-conversion}

\simulator 의 환기설비 입력은 \cref{subsubsection:ioref:ventilationsystem}과 같이 난방 및 냉방 열회수효율을 받는다. 이 값으로 존의 외기유량과 배기팬 압력상승을 계산하고, 존별 \texttt{ZoneVentilation:DesignFlowRate} 객체를 생성한다 \cite{energyplusior242}.

\subsection{입력 해석과 기본값}

환기설비는 ID로 생성된 뒤 각 존의 \inlinecode{ventilation_system_id}를 통해 연결된다. 난방 열회수효율 \inlinecode{efficiency_heating}과 냉방 열회수효율 \inlinecode{efficiency_cooling}의 기본값은 각각 0.70과 0.45이다. 두 값은 0보다 크고 1 이하인 무차원 비율로 검증된다. 환기설비가 연결되지 않은 존에도 재실 프로필이 존재하면 기준 외기유량을 갖는 존 환기 설계유량 객체가 생성된다.

\subsection{등가 외기유량과 팬동력 변환}

환기설비가 없는 존은 EnergyPlus의 \texttt{Flow/Person} 방법을 사용하고, 인당 외기유량을 다음과 같이 고정한다.

\begin{equation}
  \dot V_{\mathrm{base,person}} = 0.0083 \qquad [\mathrm{m^3/(s\cdot person)}]
  \label{eq:ventilation-base-flow}
\end{equation}

환기설비가 연결되면 난방 및 냉방 효율의 산술평균을 전체 열회수효율로 정의한다.

\begin{equation}
  \eta_{\mathrm{avg}} = \frac{\eta_{\mathrm{heating}}+\eta_{\mathrm{cooling}}}{2}
  \label{eq:ventilation-average-efficiency}
\end{equation}

EnergyPlus에 입력되는 인당 등가 외기유량과 팬 압력상승은 다음과 같다. 이 때 팬 전효율은 0.85, 환기 유형은 \texttt{Exhaust}로 고정한다. 

\begin{align}
  \dot V_{\mathrm{eq,person}} &= 0.0083\left(1-\eta_{\mathrm{avg}}\right)
  && [\mathrm{m^3/(s\cdot person)}],
  \label{eq:ventilation-equivalent-flow} \\
  \Delta P_{\mathrm{fan}} &= \frac{50}{1-\eta_{\mathrm{avg}}}
  && [\mathrm{Pa}].
  \label{eq:ventilation-fan-pressure}
\end{align}

존 전체 설계 외기유량은 EnergyPlus가 \texttt{People} 객체의 설계 인원수와 \cref{eq:ventilation-base-flow,eq:ventilation-equivalent-flow}의 인당 유량을 곱하여 계산한다. \texttt{People} 객체의 설계 인원밀도는 프로필의 재실자 스케줄 최댓값을 사용한다. 환기 객체에는 가용스케줄명을 전달하지 않으므로, \pymodule의 재실스케줄을 환기 가용스케줄로 직접 연결하지 않는다.

\subsection{EnergyPlus 객체 구성}

\begin{defaulttable}
  \caption{환기 입력에 따른 EnergyPlus 객체 필드}
  \label{tab:ventilation-object-mapping}
  \small
  \begin{tabularx}{\textwidth}{p{3.0cm}p{3.2cm}X}
    \toprule
    조건 & 필드 & 생성값 \\ \midrule
    환기설비 없음 & 계산방법 / 인당유량 & \texttt{Flow/Person} / $0.0083\,\mathrm{m^3/(s\cdot person)}$ \\
    환기설비 있음 & 계산방법 / 인당유량 & \texttt{Flow/Person} / $0.0083(1-\eta_{\mathrm{avg}})$ \\
    환기설비 있음 & 환기유형 & \texttt{Exhaust} \\
    환기설비 있음 & 팬 압력상승 & $50/(1-\eta_{\mathrm{avg}})\,\mathrm{Pa}$ \\
    환기설비 있음 & 팬 전효율 & 0.85 \\
    공통 & 존 연결 & 존별 \texttt{ZoneVentilation:DesignFlowRate} 한 개 \\ \bottomrule
  \end{tabularx}
\end{defaulttable}

\section{신재생설비}
\label{section:renewable-conversion}

현 버전에서 건물 단위로 직접 생성되는 신재생설비는 태양광발전(PV)이다. 입력 항목은 패널 면적, 변환효율, 경사각 및 방위각이며, 각 PV 입력은 독립된 발전기와 전력분배계통으로 변환된다. 지열히트펌프는 입력 유형에 포함되지만 \cref{section:hvac-conversion}에서 기술한 바와 같이 일반 전기식 VRF로 등가화되므로, 본 절의 물리적 신재생발전 모델에는 포함되지 않는다.

\subsection{PV 입력과 기하 형상 변환}

PV 면적 $A_{\mathrm{PV}}$는 정사각형 차양면의 가로와 세로 길이로 변환한다.

\begin{equation}
  L_{\mathrm{PV}}=H_{\mathrm{PV}}=\sqrt{A_{\mathrm{PV}}}.
  \label{eq:pv-square-dimension}
\end{equation}

이 면은 실제 지붕면이나 외벽면에 부착하지 않는다. 원점 기준 시작점 $(0,0,10\,\mathrm{m})$을 갖는 독립 \texttt{Shading:}\allowbreak\texttt{Site} 객체로 생성한다. 입력 \inlinecode{azimuth}와 \inlinecode{tilt}는 이 직사각형의 방위각과 경사각 필드에 직접 전달된다. 따라서 여러 PV 입력은 각각 별도의 정사각형 차양면을 생성하며, 입력만으로 건물 지붕의 위치·높이·경계와 기하학적으로 결합되지는 않는다. 건물의 전역 일사계산 설정은 \texttt{MinimalShadowing}이므로, 복잡한 상호음영보다는 단순한 발전면 일사 수집을 목적으로 한다.

\subsection{전기적 객체 구성과 발전량 계산}

단순 태양광 성능 객체(\texttt{PhotovoltaicPerformance:}\allowbreak\texttt{Simple})는 EnergyPlus의 \texttt{Fixed} 효율모드를 사용한다. 사용자 입력 \inlinecode{efficiency}는 고정 변환효율에 직접 전달되고, 전체 표면 중 활성 태양전지 면적의 비율은 0.70으로 고정된다. 순간 발전출력에는 다음 관계를 적용한다.

\begin{equation}
  P_{\mathrm{PV}} \simeq I_{\mathrm{POA}}\,A_{\mathrm{PV}}\,(0.70)\,\eta_{\mathrm{PV}},
  \label{eq:pv-simple-power}
\end{equation}

여기서 $I_{\mathrm{POA}}$는 EnergyPlus가 \texttt{Shading:Site} 면에 대해 계산한 면일사량이고, $\eta_{\mathrm{PV}}$는 입력 변환효율이다. \texttt{Generator:Photovoltaic}은 이 표면과 성능 객체를 연결한다.

각 PV에는 전용 발전기 목록, 단순 인버터 및 전력분배 객체를 만든다. 발전기 목록과 배전계통의 정격 상한은 $1{,}000{,}000\,\mathrm{W}$로 설정한다. 인버터는 항상 가동되는 \texttt{ALLON} 스케줄, 대기전력 0 W, 정격효율 1.0을 사용한다. 전력버스 유형은 \texttt{DirectCurrentWithInverter}, 발전운전방식은 \texttt{Baseload}이다. 즉, 발전 가능한 태양광 전력을 우선 부하에 공급하며, 별도의 축전지나 변압기 객체는 생성하지 않는다.

\begin{defaulttable}
  \caption{PV 입력값과 EnergyPlus 객체 필드의 대응}
  \label{tab:pv-input-mapping}
  \footnotesize
  \begin{tabularx}{\textwidth}{>{\raggedright\arraybackslash}p{2.65cm}>{\raggedright\arraybackslash}p{4.35cm}X}
    \toprule
    입력 또는 고정값 & EnergyPlus 객체 & 적용 방법 \\ \midrule
    \texttt{area} & \texttt{Shading:Site} & 가로·세로를 각각 $\sqrt{A_{\mathrm{PV}}}$로 설정 \\
    \texttt{azimuth} & \texttt{Shading:Site} & 방위각에 직접 전달 \\
    \texttt{tilt} & \texttt{Shading:Site} & 경사각에 직접 전달 \\
    \texttt{efficiency} & \makecell[l]{\texttt{Photovoltaic}\\\texttt{Performance:Simple}} & \texttt{Fixed} 변환효율에 직접 전달 \\  \bottomrule
  \end{tabularx}
\end{defaulttable}


\section{기상데이터} \label{section:weather-data}
\subsection{주소 및 행정구역 기준}
주소 정보는 2024년 12월 31일 기준 행정구역 체계를 적용하였으며, 세종특별자치시와 제주특별자치도(제주시, 서귀포시)를 포함하여 전국을 총 229개 시·군·구로 분류하였다.
각 시·군·구의 중심 좌표는 Naver Cloud Platform(NCP)에서 제공하는 Geocoding API를 기반으로 한 \texttt{korean\_geocoding} 파이썬 패키지를 활용하여 위도(latitude)와 경도(longitude) 정보를 추출하였다.

\subsection{기상데이터 수집}
기상데이터는 국제적으로 통용되는 기상데이터 플랫폼인 OneBuilding.org\cite{onebuilding2025}에서 제공하는 TMYx(Typical Meteorological Year) 데이터를 활용하였다. 해당 플랫폼은 각 지역별 기상데이터를 총 7개 형식(clm, ddy, epw, pvsyst, rain, stat, wea)으로 제공하며, 본 시뮬레이터에서는 EnergyPlus 시뮬레이션에 요구되는 EPW 형식을 사용한다.
\begin{itemize}
  \item 수집 대상: 전국 95개 기상지점
  \item 데이터 기간: 2009년 - 2023년
  \item 데이터 검토: 결측치 없음
\end{itemize}
\paragraph{tmy} TMYx(Typical Meteorological Year)는 ISO 15927-4를 기반으로 작성한 표준 기상데이터를 의미한다. ISO 15927-4 방법은 아래와 같은 절차로 이루어진다\cite{shimtmyx} (Figure \ref{fig:tmyxalgorithm})
\begin{enumerate}
  \item 지역별 데이터 수집 (시간별 건구온도, 상대습도, 수평면 전일사량, 풍속)
  \item 시간별 수집 데이터를 일 평균 데이터로 변환
  \item 하루평균 데이터를 long-term CDF(전체 수집 기간), monthly CDF로 구분하고, 건구온도, 상대습도, 수평면 전일사량의 누적분포 함수 계산 (FS방법)
  
\[ 
  FS_{x}(y,m) = \left(\frac{1}{n}\right) \sum_{i=1}^{n} |CDF_{m}(x_{i}) - CDF_{y,m}(x_{i})|
\]

  \item $FS_{x}(y,m)$의 값이 작은 순으로 월 별 3개 연도를 선정
  \item 장기 평균 풍속과 월별 평균 풍속의 편차가 가장 작은 해를 대표 해로 선정
\end{enumerate}

\begin{defaultfigure}
  \includegraphics[width=\textwidth]{표준기상데이터생성.png}
  \caption{표준기상데이터 생성 방법 (ISO 15927-4)}
  \label{fig:tmyxalgorithm}
\end{defaultfigure}


\subsection{주소와 기상데이터 매칭방식}
\label{sec:address_weather_matching}
기상데이터 매칭은 각 시·군·구의 중심 좌표와 95개 TMY 기상지점의 위도·경도 정보를 비교하여, Haversine 공식을 이용한 거리 계산 결과 가장 가까운 지점을 선택하는 방식으로 매칭함. 
각 행정구역별 최근접 기상지점의 epw를 매칭한 결과는 표~\ref{tab:weather_mapping}에 제시함.

\renewcommand{\arraystretch}{0.8}
\begin{longtable}{lcclcc}
  \caption{행정구역별 기상데이터 매칭 결과} \\
  \label{tab:weather_mapping} \\
  \toprule
  \small
  시·군·구 & 위도[$^\circ$] & 경도[$^\circ$] & 매칭 EPW 지역명 & WMOcode & EPW 기준년도 \\ \midrule
  \endfirsthead
  \multicolumn{6}{r}{\textit{(이전 페이지에서 계속)}} \\ \toprule
  시·군·구 & 위도[$^\circ$] & 경도[$^\circ$] & 매칭 EPW 지역명 & WMOcode & EPW 기준년도 \\ \midrule
  \endhead
  \midrule \multicolumn{6}{r}{\textit{(다음 페이지에 계속)}} \\ \bottomrule
  \endfoot
  \bottomrule
  \endlastfoot
  % 본문 데이터 (예시)
  서울특별시 종로구 & 37.57 & 126.98 & Seoul.WS & 471080 & '09-'23 \\
  서울특별시 중구 & 37.56 & 127.00 & Seoul.WS & 471080 & '09-'23 \\
  서울특별시 용산구 & 37.53 & 126.99 & Seoul.WS & 471080 & '09-'23 \\
  서울특별시 성동구 & 37.56 & 127.04 & Seoul.WS & 471080 & '09-'23 \\
  서울특별시 광진구 & 37.54 & 127.08 & Seoul-Seongnam.AP & 471110 & '09-'23 \\
  서울특별시 동대문구 & 37.57 & 127.04 & Seoul.WS & 471080 & '09-'23 \\
  서울특별시 중랑구 & 37.61 & 127.09 & Seoul.WS & 471080 & '09-'23 \\
  서울특별시 성북구 & 37.59 & 127.02 & Seoul.WS & 471080 & '09-'23 \\
  서울특별시 강북구 & 37.64 & 127.03 & Seoul.WS & 471080 & '09-'23 \\
  서울특별시 도봉구 & 37.67 & 127.05 & Seoul.WS & 471080 & '09-'23 \\
  서울특별시 노원구 & 37.65 & 127.06 & Seoul.WS & 471080 & '09-'23 \\
  서울특별시 은평구 & 37.60 & 126.93 & Seoul.WS & 471080 & '09-'23 \\
  서울특별시 서대문구 & 37.58 & 126.94 & Seoul.WS & 471080 & '09-'23 \\
  서울특별시 마포구 & 37.57 & 126.90 & Seoul.WS & 471080 & '09-'23 \\
  서울특별시 양천구 & 37.52 & 126.87 & Gimpo.Intl.AP & 471100 & '09-'23 \\
  서울특별시 강서구 & 37.55 & 126.85 & Gimpo.Intl.AP & 471100 & '09-'23 \\
  서울특별시 구로구 & 37.50 & 126.89 & Seoul.WS & 471080 & '09-'23 \\
  서울특별시 금천구 & 37.46 & 126.90 & Seoul.WS & 471080 & '09-'23 \\
  서울특별시 영등포구 & 37.53 & 126.90 & Seoul.WS & 471080 & '09-'23 \\
  서울특별시 동작구 & 37.51 & 126.94 & Seoul.WS & 471080 & '09-'23 \\
  서울특별시 관악구 & 37.48 & 126.95 & Seoul.WS & 471080 & '09-'23 \\
  서울특별시 서초구 & 37.48 & 127.03 & Seoul-Seongnam.AP & 471110 & '09-'23 \\
  서울특별시 강남구 & 37.52 & 127.05 & Seoul.WS & 471080 & '09-'23 \\
  서울특별시 송파구 & 37.51 & 127.11 & Seoul-Seongnam.AP & 471110 & '09-'23 \\
  서울특별시 강동구 & 37.53 & 127.12 & Seoul-Seongnam.AP & 471110 & '09-'23 \\
  부산광역시 중구 & 35.11 & 129.03 & Busan-Daecheongdong.WS & 471590 & '09-'23 \\
  부산광역시 서구 & 35.10 & 129.02 & Busan-Daecheongdong.WS & 471590 & '09-'23 \\
  부산광역시 동구 & 35.13 & 129.05 & Busan-Daecheongdong.WS & 471590 & '09-'23 \\
  부산광역시 영도구 & 35.09 & 129.07 & Busan-Daecheongdong.WS & 471590 & '09-'23 \\
  부산광역시 부산진구 & 35.16 & 129.05 & Busan-Daecheongdong.WS & 471590 & '09-'23 \\
  부산광역시 동래구 & 35.20 & 129.09 & Busan-Daecheongdong.WS & 471590 & '09-'23 \\
  부산광역시 남구 & 35.14 & 129.08 & Busan-Daecheongdong.WS & 471590 & '09-'23 \\
  부산광역시 북구 & 35.20 & 128.99 & Busan-Gimhae.Intl.AP & 471530 & '09-'23 \\
  부산광역시 해운대구 & 35.16 & 129.16 & Busan-Daecheongdong.WS & 471590 & '09-'23 \\
  부산광역시 사하구 & 35.10 & 128.97 & Busan-Daecheongdong.WS & 471590 & '09-'23 \\
  부산광역시 금정구 & 35.24 & 129.09 & Busan-Gimhae.Intl.AP & 471530 & '09-'23 \\
  부산광역시 강서구 & 35.21 & 128.98 & Busan-Gimhae.Intl.AP & 471530 & '09-'23 \\
  부산광역시 연제구 & 35.18 & 129.08 & Busan-Daecheongdong.WS & 471590 & '09-'23 \\
  부산광역시 수영구 & 35.15 & 129.11 & Busan-Daecheongdong.WS & 471590 & '09-'23 \\
  부산광역시 사상구 & 35.15 & 128.99 & Busan-Gimhae.Intl.AP & 471530 & '09-'23 \\
  부산광역시 기장군 & 35.24 & 129.22 & Busan-Daecheongdong.WS & 471590 & '09-'23 \\
  대구광역시 중구 & 35.87 & 128.61 & Camp.Walker-Daegu & 471425 & '09-'23 \\
  대구광역시 동구 & 35.89 & 128.64 & Daegu.Intl.AP & 471420 & '09-'23 \\
  대구광역시 서구 & 35.87 & 128.56 & Camp.Walker-Daegu & 471425 & '09-'23 \\
  대구광역시 남구 & 35.85 & 128.60 & Camp.Walker-Daegu & 471425 & '09-'23 \\
  대구광역시 북구 & 35.89 & 128.58 & Camp.Walker-Daegu & 471425 & '09-'23 \\
  대구광역시 수성구 & 35.86 & 128.63 & Daegu & 471430 & '09-'23 \\
  대구광역시 달서구 & 35.83 & 128.53 & Camp.Walker-Daegu & 471425 & '09-'23 \\
  대구광역시 달성군 & 35.77 & 128.43 & Camp.Walker-Daegu & 471425 & '09-'23 \\
  대구광역시 군위군 & 36.24 & 128.57 & Gumi & 471540 & '09-'23 \\
  인천광역시 강화군 & 37.75 & 126.49 & Pyoripsan & 471084 & '09-'23 \\
  인천광역시 옹진군 & 37.45 & 126.64 & Incheon.WS & 471120 & '09-'23 \\
  인천광역시 중구 & 37.47 & 126.62 & Incheon.WS & 471120 & '09-'23 \\
  인천광역시 동구 & 37.47 & 126.64 & Incheon.WS & 471120 & '09-'23 \\
  인천광역시 미추홀구 & 37.46 & 126.65 & Incheon.WS & 471120 & '09-'23 \\
  인천광역시 연수구 & 37.41 & 126.68 & Incheon.WS & 471120 & '09-'23 \\
  인천광역시 남동구 & 37.45 & 126.73 & Incheon.WS & 471120 & '09-'23 \\
  인천광역시 부평구 & 37.51 & 126.72 & Gimpo.Intl.AP & 471100 & '09-'23 \\
  인천광역시 계양구 & 37.54 & 126.74 & Gimpo.Intl.AP & 471100 & '09-'23 \\
  인천광역시 서구 & 37.55 & 126.68 & Incheon.WS & 471120 & '09-'23 \\
  광주광역시 동구 & 35.15 & 126.92 & Gwangju & 471560 & '09-'23 \\
  광주광역시 서구 & 35.15 & 126.89 & Gwangju & 471560 & '09-'23 \\
  광주광역시 남구 & 35.13 & 126.90 & Gwangju & 471560 & '09-'23 \\
  광주광역시 북구 & 35.17 & 126.91 & Gwangju & 471560 & '09-'23 \\
  광주광역시 광산구 & 35.14 & 126.79 & Gwangju.AP & 471580 & '09-'23 \\
  대전광역시 동구 & 36.31 & 127.45 & Daejeon.WS & 471330 & '09-'23 \\
  대전광역시 중구 & 36.33 & 127.42 & Daejeon.WS & 471330 & '09-'23 \\
  대전광역시 서구 & 36.36 & 127.38 & Daejeon.WS & 471330 & '09-'23 \\
  대전광역시 유성구 & 36.36 & 127.36 & Daejeon.WS & 471330 & '09-'23 \\
  대전광역시 대덕구 & 36.35 & 127.42 & Daejeon.WS & 471330 & '09-'23 \\
  울산광역시 중구 & 35.57 & 129.33 & Ulsan & 471520 & '09-'23 \\
  울산광역시 남구 & 35.54 & 129.33 & Ulsan & 471520 & '09-'23 \\
  울산광역시 동구 & 35.50 & 129.42 & Ulsan & 471520 & '09-'23 \\
  울산광역시 북구 & 35.58 & 129.36 & Ulsan & 471520 & '09-'23 \\
  울산광역시 울주군 & 35.52 & 129.24 & Ulsan & 471520 & '09-'23 \\
  세종특별자치시 & 36.48 & 127.29 & Gongju & 692214 & '09-'23 \\
  경기도 수원시 & 37.26 & 127.03 & Suwon.AP & 471200 & '09-'23 \\
  경기도 고양시 & 37.66 & 126.83 & Gimpo.Intl.AP & 471100 & '09-'23 \\
  경기도 용인시 & 37.24 & 127.18 & Maesanri & 471205 & '09-'23 \\
  경기도 성남시 & 37.42 & 127.13 & Seoul-Seongnam.AP & 471110 & '09-'23 \\
  경기도 부천시 & 37.50 & 126.77 & Gimpo.Intl.AP & 471100 & '09-'23 \\
  경기도 화성시 & 37.20 & 126.83 & Suwon.WS & 471190 & '09-'23 \\
  경기도 안산시 & 37.32 & 126.83 & Suwon.WS & 471190 & '09-'23 \\
  경기도 남양주시 & 37.64 & 127.22 & Yangsu.Ri & 693594 & '09-'23 \\
  경기도 안양시 & 37.39 & 126.96 & Suwon.WS & 471190 & '09-'23 \\
  경기도 평택시 & 36.99 & 127.11 & Pyeongtaek.AP & 471270 & '09-'23 \\
  경기도 시흥시 & 37.38 & 126.80 & Incheon.WS & 471120 & '09-'23 \\
  경기도 파주시 & 37.76 & 126.78 & Paju & 470990 & '09-'23 \\
  경기도 의정부시 & 37.74 & 127.03 & Gwangjeok & 693604 & '09-'23 \\
  경기도 김포시 & 37.62 & 126.72 & Gimpo.Intl.AP & 471100 & '09-'23 \\
  경기도 광주시 & 37.43 & 127.26 & Maesanri & 471205 & '09-'23 \\
  경기도 광명시 & 37.48 & 126.86 & Gimpo.Intl.AP & 471100 & '09-'23 \\
  경기도 군포시 & 37.36 & 126.94 & Suwon.WS & 471190 & '09-'23 \\
  경기도 하남시 & 37.54 & 127.21 & Seoul-Seongnam.AP & 471110 & '09-'23 \\
  경기도 오산시 & 37.15 & 127.08 & Osan.AB & 471220 & '09-'23 \\
  경기도 양주시 & 37.79 & 127.05 & Gwangjeok & 693604 & '09-'23 \\
  경기도 이천시 & 37.27 & 127.44 & Icheon & 470970 & '09-'23 \\
  경기도 구리시 & 37.59 & 127.13 & Seoul.WS & 471080 & '09-'23 \\
  경기도 안성시 & 37.01 & 127.28 & Pyeongtaek.AP & 471270 & '09-'23 \\
  경기도 포천시 & 37.89 & 127.20 & Dongducheon & 470980 & '09-'23 \\
  경기도 의왕시 & 37.34 & 126.97 & Suwon.WS & 471190 & '09-'23 \\
  경기도 양평군 & 37.49 & 127.49 & Yeoju.Range & 471315 & '09-'23 \\
  경기도 여주시 & 37.30 & 127.60 & Icheon & 470970 & '09-'23 \\
  경기도 동두천시 & 37.90 & 127.06 & Dongducheon & 470980 & '09-'23 \\
  경기도 과천시 & 37.43 & 126.99 & Seoul-Seongnam.AP & 471110 & '09-'23 \\
  경기도 가평군 & 37.83 & 127.51 & Gapyeong & 692114 & '09-'23 \\
  경기도 연천군 & 38.10 & 127.07 & Cheorwon & 470950 & '09-'23 \\
  강원도 춘천시 & 37.88 & 127.73 & Chuncheon-Camp.Page.AP & 471040 & '09-'23 \\
  강원도 원주시 & 37.34 & 127.92 & Wonju.WS & 471140 & '09-'23 \\
  강원도 강릉시 & 37.75 & 128.88 & Gangneung & 471050 & '09-'23 \\
  강원도 동해시 & 37.52 & 129.11 & Donghae.Radar & 471060 & '09-'23 \\
  강원도 태백시 & 37.16 & 128.99 & Taebaek.Provincial.Park & 471223 & '09-'23 \\
  강원도 속초시 & 38.21 & 128.59 & Sokcho & 470900 & '09-'23 \\
  강원도 삼척시 & 37.45 & 129.17 & Donghae.Radar & 471060 & '09-'23 \\
  강원도 홍천군 & 37.70 & 127.89 & Hongchon & 693624 & '09-'23 \\
  강원도 횡성군 & 37.49 & 127.99 & Wonju.WS & 471140 & '09-'23 \\
  강원도 영월군 & 37.18 & 128.46 & Yeongwol.WS & 471210 & '09-'23 \\
  강원도 평창군 & 37.37 & 128.39 & Yeongwol.WS & 471210 & '09-'23 \\
  강원도 정선군 & 37.38 & 128.66 & Yeongwol.WS & 471210 & '09-'23 \\
  강원도 철원군 & 38.15 & 127.31 & Cheorwon & 470950 & '09-'23 \\
  강원도 화천군 & 38.11 & 127.71 & Sachang.Ri & 693614 & '09-'23 \\
  강원도 양구군 & 38.11 & 127.99 & Yanggu & 692094 & '09-'23 \\
  강원도 인제군 & 38.07 & 128.17 & Yanggu & 692094 & '09-'23 \\
  강원도 고성군 & 38.38 & 128.47 & Geojin & 470800 & '09-'23 \\
  강원도 양양군 & 38.08 & 128.62 & Yangyang.Intl.AP & 470920 & '09-'23 \\
  충청북도 청주시 & 36.64 & 127.49 & Cheongju.WS & 471310 & '09-'23 \\
  충청북도 충주시 & 36.99 & 127.93 & Jungwon.AB & 471250 & '09-'23 \\
  충청북도 제천시 & 37.13 & 128.19 & Yeongwol.WS & 471210 & '09-'23 \\
  충청북도 보은군 & 36.49 & 127.73 & Seongmu.AB & 471240 & '09-'23 \\
  충청북도 옥천군 & 36.31 & 127.57 & Daejeon.WS & 471330 & '09-'23 \\
  충청북도 영동군 & 36.18 & 127.78 & Chupungnyeong & 471350 & '09-'23 \\
  충청북도 증평군 & 36.79 & 127.58 & Cheongju.Intl.AP & 471280 & '09-'23 \\
  충청북도 진천군 & 36.86 & 127.44 & Cheonan & 471450 & '09-'23 \\
  충청북도 괴산군 & 36.82 & 127.79 & Jungwon.AB & 471250 & '09-'23 \\
  충청북도 음성군 & 36.94 & 127.69 & Jungwon.AB & 471250 & '09-'23 \\
  충청북도 단양군 & 36.98 & 128.37 & Yeongwol.WS & 471210 & '09-'23 \\
  충청남도 천안시 & 36.82 & 127.11 & Cheonan & 471450 & '09-'23 \\
  충청남도 공주시 & 36.45 & 127.12 & Gongju & 692214 & '09-'23 \\
  충청남도 보령시 & 36.33 & 126.61 & Boryeong & 471500 & '09-'23 \\
  충청남도 아산시 & 36.79 & 127.00 & Pyeongtaek.AP & 471270 & '09-'23 \\
  충청남도 서산시 & 36.78 & 126.45 & Seosan.WS & 471290 & '09-'23 \\
  충청남도 논산시 & 36.19 & 127.10 & Daejeon.WS & 471330 & '09-'23 \\
  충청남도 계룡시 & 36.27 & 127.25 & Daejeon.WS & 471330 & '09-'23 \\
  충청남도 당진시 & 36.89 & 126.65 & Seosan.WS & 471290 & '09-'23 \\
  충청남도 금산군 & 36.11 & 127.49 & Geumsan & 692244 & '09-'23 \\
  충청남도 부여군 & 36.28 & 126.91 & Boryeong & 471500 & '09-'23 \\
  충청남도 서천군 & 36.08 & 126.69 & Gunsan & 471400 & '09-'23 \\
  충청남도 청양군 & 36.46 & 126.80 & Boryeong & 471500 & '09-'23 \\
  충청남도 홍성군 & 36.60 & 126.66 & Seosan.WS & 471290 & '09-'23 \\
  충청남도 예산군 & 36.68 & 126.84 & Seosan.WS & 471290 & '09-'23 \\
  충청남도 태안군 & 36.75 & 126.30 & Seosan.WS & 471290 & '09-'23 \\
  전라북도 전주시 & 35.82 & 127.15 & Jeonju & 471460 & '09-'23 \\
  전라북도 군산시 & 35.97 & 126.74 & Gunsan & 471400 & '09-'23 \\
  전라북도 익산시 & 35.95 & 126.96 & Jeonju & 471460 & '09-'23 \\
  전라북도 정읍시 & 35.57 & 126.86 & Jeongeup & 471710 & '09-'23 \\
  전라북도 남원시 & 35.42 & 127.39 & Namwon & 471730 & '09-'23 \\
  전라북도 김제시 & 35.80 & 126.88 & Jeonju & 471460 & '09-'23 \\
  전라북도 완주군 & 35.90 & 127.16 & Jeonju & 471460 & '09-'23 \\
  전라북도 진안군 & 35.79 & 127.42 & Jeonju & 471460 & '09-'23 \\
  전라북도 무주군 & 36.01 & 127.66 & Geumsan & 692244 & '09-'23 \\
  전라북도 장수군 & 35.65 & 127.52 & Namwon & 471730 & '09-'23 \\
  전라북도 임실군 & 35.62 & 127.29 & Namwon & 471730 & '09-'23 \\
  전라북도 순창군 & 35.37 & 127.14 & Namwon & 471730 & '09-'23 \\
  전라북도 고창군 & 35.44 & 126.70 & Gochang & 471720 & '09-'23 \\
  전라북도 부안군 & 35.73 & 126.73 & Jeongeup & 471710 & '09-'23 \\
  전라남도 목포시 & 34.81 & 126.39 & Mokpo & 471650 & '09-'23 \\
  전라남도 여수시 & 34.76 & 127.66 & Yeosu.Obs & 471680 & '09-'23 \\
  전라남도 순천시 & 34.95 & 127.49 & Suncheon.WS & 471740 & '09-'23 \\
  전라남도 나주시 & 35.02 & 126.71 & Gwangju.AP & 471580 & '09-'23 \\
  전라남도 광양시 & 34.94 & 127.70 & Yeosu.AP & 471670 & '09-'23 \\
  전라남도 담양군 & 35.32 & 126.99 & Gwangju & 471560 & '09-'23 \\
  전라남도 곡성군 & 35.28 & 127.29 & Namwon & 471730 & '09-'23 \\
  전라남도 구례군 & 35.20 & 127.46 & Suncheon.WS & 471740 & '09-'23 \\
  전라남도 고흥군 & 34.60 & 127.28 & Yeosu.AP & 471670 & '09-'23 \\
  전라남도 보성군 & 34.77 & 127.08 & Suncheon.WS & 471740 & '09-'23 \\
  전라남도 화순군 & 35.06 & 126.99 & Gwangju & 471560 & '09-'23 \\
  전라남도 장흥군 & 34.68 & 126.91 & Wando.WS & 471700 & '09-'23 \\
  전라남도 강진군 & 34.64 & 126.77 & Wando.WS & 471700 & '09-'23 \\
  전라남도 해남군 & 34.57 & 126.60 & Wando.WS & 471700 & '09-'23 \\
  전라남도 영암군 & 34.80 & 126.70 & Mokpo & 471650 & '09-'23 \\
  전라남도 무안군 & 34.99 & 126.48 & Muan.Intl.AP & 471630 & '09-'23 \\
  전라남도 함평군 & 35.07 & 126.52 & Muan.Intl.AP & 471630 & '09-'23 \\
  전라남도 영광군 & 35.28 & 126.51 & Gochang & 471720 & '09-'23 \\
  전라남도 장성군 & 35.30 & 126.78 & Gwangju & 471560 & '09-'23 \\
  전라남도 완도군 & 34.31 & 126.75 & Wando.WS & 471700 & '09-'23 \\
  전라남도 진도군 & 34.49 & 126.26 & Jindo.WS & 471750 & '09-'23 \\
  전라남도 신안군 & 34.83 & 126.35 & Mokpo & 471650 & '09-'23 \\
  경상북도 포항시 & 36.02 & 129.34 & Pohang.WS & 471380 & '09-'23 \\
  경상북도 경주시 & 35.86 & 129.22 & Pohang.AP & 471390 & '09-'23 \\
  경상북도 김천시 & 36.14 & 128.11 & Chupungnyeong & 471350 & '09-'23 \\
  경상북도 안동시 & 36.57 & 128.73 & Andong.WS & 471360 & '09-'23 \\
  경상북도 구미시 & 36.12 & 128.34 & Gumi & 471540 & '09-'23 \\
  경상북도 영주시 & 36.81 & 128.62 & Andong.WS & 471360 & '09-'23 \\
  경상북도 영천시 & 35.97 & 128.94 & Daegu.Intl.AP & 471420 & '09-'23 \\
  경상북도 상주시 & 36.41 & 128.16 & Sangju.WS & 471370 & '09-'23 \\
  경상북도 문경시 & 36.59 & 128.19 & Yecheon.AP & 471340 & '09-'23 \\
  경상북도 경산시 & 35.83 & 128.74 & Daegu & 471430 & '09-'23 \\
  경상북도 의성군 & 36.35 & 128.70 & Andong.WS & 471360 & '09-'23 \\
  경상북도 청송군 & 36.44 & 129.06 & Andong.WS & 471360 & '09-'23 \\
  경상북도 영양군 & 36.67 & 129.11 & Andong.WS & 471360 & '09-'23 \\
  경상북도 영덕군 & 36.42 & 129.37 & Pohang.WS & 471380 & '09-'23 \\
  경상북도 청도군 & 35.65 & 128.73 & Cheongdo & 693414 & '09-'23 \\
  경상북도 고령군 & 35.73 & 128.26 & Camp.Walker-Daegu & 471425 & '09-'23 \\
  경상북도 성주군 & 35.92 & 128.28 & Gumi & 471540 & '09-'23 \\
  경상북도 칠곡군 & 36.00 & 128.40 & Gumi & 471540 & '09-'23 \\
  경상북도 예천군 & 36.65 & 128.44 & Yecheon.AP & 471340 & '09-'23 \\
  경상북도 봉화군 & 36.89 & 128.73 & Taebaek.Provincial.Park & 471223 & '09-'23 \\
  경상북도 울진군 & 36.99 & 129.40 & Uljin.WS & 471300 & '09-'23 \\
  경상북도 울릉군 & 37.48 & 130.91 & Ulleungdo.WS & 471150 & '09-'23 \\
  경상남도 창원시 & 35.23 & 128.68 & Changwon & 471550 & '09-'23 \\
  경상남도 진주시 & 35.18 & 128.11 & Jinju.WS & 471920 & '09-'23 \\
  경상남도 통영시 & 34.85 & 128.43 & Tongyeong.WS & 471620 & '09-'23 \\
  경상남도 사천시 & 35.00 & 128.06 & Sacheon.AB & 471610 & '09-'23 \\
  경상남도 김해시 & 35.23 & 128.89 & Busan-Gimhae.Intl.AP & 471530 & '09-'23 \\
  경상남도 밀양시 & 35.50 & 128.75 & Cheongdo & 693414 & '09-'23 \\
  경상남도 거제시 & 34.88 & 128.62 & Tongyeong.WS & 471620 & '09-'23 \\
  경상남도 양산시 & 35.34 & 129.04 & Busan-Gimhae.Intl.AP & 471530 & '09-'23 \\
  경상남도 의령군 & 35.32 & 128.26 & Jinju.WS & 471920 & '09-'23 \\
  경상남도 함안군 & 35.27 & 128.41 & Changwon & 471550 & '09-'23 \\
  경상남도 창녕군 & 35.54 & 128.49 & Cheongdo & 693414 & '09-'23 \\
  경상남도 고성군 & 34.97 & 128.32 & Tongyeong.WS & 471620 & '09-'23 \\
  경상남도 남해군 & 34.84 & 127.89 & Yeosu.Obs & 471680 & '09-'23 \\
  경상남도 하동군 & 35.07 & 127.75 & Jinju.WS & 471920 & '09-'23 \\
  경상남도 산청군 & 35.42 & 127.87 & Geochang.WS & 471570 & '09-'23 \\
  경상남도 함양군 & 35.52 & 127.73 & Geochang.WS & 471570 & '09-'23 \\
  경상남도 거창군 & 35.69 & 127.91 & Geochang.WS & 471570 & '09-'23 \\
  경상남도 합천군 & 35.57 & 128.17 & Geochang.WS & 471570 & '09-'23 \\
  제주특별자치도 제주시 & 33.50 & 126.53 & Jeju.WS & 471840 & '09-'23 \\
  제주특별자치도 서귀포시 & 33.25 & 126.56 & Seogwipo & 471890 & '09-'23 \\
\end{longtable}



\chapter{시뮬레이션 결과의 해석}

\simulator\은 계산된 용도별, 열원별, 월별 에너지소요량을 바탕으로, 1차에너지소요량, 온실가스 배출량 및 에너지 비용을 산출하여 제시한다. 각 지표에 대한 변환 계수는 본 엔진에서 지원하는 4가지 연료(전기, 천연가스, 등유, 지역난방)에 대하여 각각 정의되며, 건물 전체의 1차에너지소요량, 온실가스배출량 및 에너지 비용은 4가지 연료에 대한 값의 합으로 정의된다 (식 \ref{eq:1st-energy-calculation} - \ref{eq:cost-calculation}).

\begin{align}
E_{\mathrm{area}} &= \frac{\sum_{f} E_f \, c_f}{A} \label{eq:1st-energy-calculation} \\
G_{\mathrm{area}} &= \frac{\sum_{f} E_f \, g_f}{A} \label{eq:ghg-calculation} \\
C_{\mathrm{area}} &= \frac{\sum_{f} E_f \, p_f}{A} \label{eq:cost-calculation}
\end{align}

여기서,

\begin{description}
\item[$E_{\mathrm{area}}$] 건물 면적당 1차 에너지소요량
\item[$G_{\mathrm{area}}$] 건물 면적당 온실가스 배출량
\item[$C_{\mathrm{area}}$] 건물 면적당 에너지 비용
\item[$E_f$] 연료 $f$의 에너지 사용량
\item[$c_f$] 연료 $f$의 1차 에너지 변환계수
\item[$g_f$] 연료 $f$의 온실가스 배출계수
\item[$p_f$] 연료 $f$의 단위 에너지 비용
\item[$A$] 건물의 바닥면적
\item[$f$] 연료 유형 (전력, 가스, 지역난방, 유류)
\end{description}

\subsection{바닥면적} \label{subsec:floorarea_definition_for_EUI}

바닥면적은 각 실에 속한 면 중, `바닥(floor)'으로 정의된 면적의 합이다. 각 실의 냉난방 여부와는 무관하게, 입력된 모든 실의 바닥 면적을 합하여 계산한다.

\subsection{1차에너지소요량}
연료별 1차에너지 환산 계수는 표 \ref{tbl:coeff_kWh_to_1stkWh}\와 같다. 이는 제로에너지건축물 인증제도 운영규정 \cite{law:energy_to_1stenergy} 에서 제시하고 있다. 국내에서는 모든 연료에 대하여 대표적인 환산계수만을 제시하고 있어 \cite{Jeong2021energydeungkeup}, 가스 및 등유에 대하여 동일한 환산계수를 적용하였다.

\begin{table}[ht]
  \caption{연료별 사용량의 1차에너지 환산 계수}
  \label{tbl:coeff_kWh_to_1stkWh}  
  \centering
  \begin{tabular}{cr}
    \toprule
    연료명 & \makecell{환산계수 \\ $[kWh/kWh]$} \\ \midrule
    전기 & 2.75 \\
    가스 & 1.1 \\
    등유 & 1.1 \\
    지역난방 & 0.728 \\ \bottomrule
  \end{tabular}
\end{table}

\subsection{온실가스 배출량}

연료별 온실가스 배출량 환산 계수는 표 \ref{tbl:coeff_kWh_to_CO2}\와 같다. 전기의 환산 계수는 환경부 온실가스정보센터에서 공시한 바에 따른다. 가스 및 등유의 환산 계수는 기후변화에 관한 정부간 패널 (Intergovernmental Panel on Climate Change; IPCC)에서 공시한 바에 따르며, 지역난방의 환산계수는 한국지역난방공사에서 공포하는 바에 따른다.

\begin{table}[ht]
  \caption{연료별 사용량의 온실가스 배출량 환산 계수}
  \label{tbl:coeff_kWh_to_CO2}  
  \centering
  \begin{tabular}{cr}
    \toprule
    연료명 & 환산계수 [kgCO\textsubscript{2,eq}/kWh] \\ \midrule
    전기 & 0.4541 \\
    가스 & 0.2024 \\
    등유 & 0.2603 \\
    지역난방 & 0.1358 \\ \bottomrule
  \end{tabular}
\end{table}

\subsection{비용}
연료별 비용 환산 계수는 표 \ref{tbl:coeff_kWh_to_cost}\와 같다. 에너지 요금 체계는 기본요금, 구간별 단가, 계절별 차등 등 다양한 요소로 구성되어 있어 단일 기준으로 정확히 산정하기는 어렵다. 따라서 복잡한 체계를 단순화하여 대표 환산계수를 적용하는 통계적 접근 방식을 사용하였다.

\begin{table}[ht]
  \caption{연료별 사용량의 요금 환산 계수}
  \label{tbl:coeff_kWh_to_cost}  
  \centering
  \begin{tabular}{cr}
    \toprule
    연료명 & 환산계수 [won/kWh] \\ \midrule
    전기 & 162.92 \\
    가스 & 78.12 \\
    등유 & 141.92 \\
    지역난방 & 94.98 \\ \bottomrule
  \end{tabular}
\end{table}

\paragraph{전기, 가스 및 등유} 매월 발간되는 KESIS 국가에너지통계 종합정보시스템에서 2024년 기준값을 참고하였다 (그림 \ref{fig:naturalgas_cost}). 전기와 가스의 경우 말단 용도별로 상이한 통계치가 제시되는데, 각 건물의 계약전력 등 세부 구분을 명확히 반영하기 어려우므로 각 용도별 값의 평균을 사용하였다.

\begin{defaultfigure}
  \includegraphics[width=\textwidth]{KIESIS 도시가스 가격 공표 (나중에 표로 바꾸기).png}
  \caption{KIESIS의 도시가스 가격 공포 (2024년)}
  \label{fig:naturalgas_cost}
\end{defaultfigure}

\paragraph{지역난방} 지역난방은 한국지역난방공사의 경영실적 공시에 명시된 판매단가에 따랐다.